% !TEX program = lualatex
% =====================================================================
%  Trabajo de Innovación — UNIR
%  Cerebro Predictivo y Optimización Logística para Máquinas
%  Expendedoras de Bebidas
%
%  Compilar: latexmk -lualatex main.tex
% =====================================================================
\documentclass[12pt,a4paper]{article}

% ---- Datos del trabajo ------------------------------------------
\newcommand{\tituloproyecto}{Cerebro Predictivo y Optimización Logística para Máquinas Expendedoras de Bebidas}
\newcommand{\autorescorto}{Camacho, Diaz, Hernández y Sais}
\newcommand{\fechaproyecto}{\today}

% ---------------------------------------------------------------
%  Preámbulo — configuración general del documento
% ---------------------------------------------------------------

% ---- Idioma y codificación ------------------------------------
\usepackage[utf8]{inputenc}
\usepackage[T1]{fontenc}
\usepackage[spanish,es-tabla,es-nodecimaldot]{babel}
\usepackage{csquotes}

% ---- Tipografía -----------------------------------------------
% El documento original usa Calibri. "carlito" es la fuente libre
% métricamente compatible; si no está instalada se usa la de por defecto.
\IfFileExists{carlito.sty}{\usepackage{carlito}\renewcommand{\familydefault}{\sfdefault}}{}
\usepackage[protrusion=true,expansion=false]{microtype}
\usepackage{amssymb}      % casillas del cuestionario ($\square$)

% ---- Página y espaciado (idénticos al .docx) -------------------
% A4, márgenes: sup. 2,5 cm · inf. 2,5 cm · izq. 3 cm · der. 2 cm
\usepackage[a4paper,top=2.5cm,bottom=2.5cm,left=3cm,right=2cm]{geometry}
\usepackage{setspace}
\onehalfspacing                       % interlineado 1,5
\setlength{\parindent}{0pt}
\setlength{\parskip}{6pt}

% ---- Color corporativo UNIR ------------------------------------
\usepackage{xcolor}
\definecolor{unirazul}{HTML}{0098CD}

% ---- Títulos numerados y en azul --------------------------------
\usepackage{titlesec}
\titleformat{\section}{\normalfont\Large\bfseries\color{unirazul}}{\thesection.}{0.5em}{}
\titleformat{\subsection}{\normalfont\large\bfseries\color{unirazul}}{\thesubsection.}{0.5em}{}
\titleformat{\subsubsection}{\normalfont\normalsize\bfseries\color{unirazul}}{\thesubsubsection.}{0.5em}{}
\titlespacing*{\section}{0pt}{18pt}{10pt}
\titlespacing*{\subsection}{0pt}{14pt}{8pt}
\titlespacing*{\subsubsection}{0pt}{12pt}{6pt}

% ---- Índice ----------------------------------------------------
\usepackage{tocloft}
\setlength{\cftsecnumwidth}{2.2em}   % hueco del número de sección
\setcounter{tocdepth}{3}
\setcounter{secnumdepth}{3}

% ---- Figuras, tablas y listas ----------------------------------
\usepackage{graphicx}
\graphicspath{{imagenes/}}
\usepackage{booktabs}
\usepackage{tabularx}
\usepackage{longtable}
\usepackage{array}
\usepackage{calc}          % anchos de columna calculados (tablas de pandoc)
\usepackage{float}
\usepackage{multirow}
\usepackage[font=small,labelfont=bf]{caption}
\usepackage{enumitem}

% Pie de tabla/figura fuera de flotante y línea de fuente
\captionsetup{justification=raggedright,singlelinecheck=false}
\newcommand{\fuente}[1]{{\par\raggedright\footnotesize\itshape #1\par}\medskip}
\setlist{itemsep=2pt,topsep=4pt}

% ---- Encabezado y pie ------------------------------------------
\usepackage{fancyhdr}
\pagestyle{fancy}
\fancyhf{}
\fancyfoot[L]{\footnotesize\color{gray}\autorescorto}
\fancyfoot[R]{\footnotesize\thepage}
\renewcommand{\headrulewidth}{0pt}
\renewcommand{\footrulewidth}{0.4pt}

% ---------------------------------------------------------------
%  BIBLIOGRAFÍA
%  Deja \overleaftrue para compilar en Overleaf con estilo APA 7
%  (biblatex-apa + biber). Cambia a \overleaffalse solo si tu
%  instalación local no tiene biblatex (usa natbib + apalike).
% ---------------------------------------------------------------
\newif\ifoverleaf
\overleaftrue

\ifoverleaf
  \usepackage[style=apa,natbib=true,backend=biber,language=spanish]{biblatex}
  \DeclareLanguageMapping{spanish}{spanish-apa}
  \addbibresource{bibliografia/referencias.bib}
  \newcommand{\imprimirbibliografia}{%
    \printbibliography[title={Referencias bibliográficas},heading=bibnumbered]}
\else
  \usepackage{natbib}
  \bibliographystyle{apalike}
  \setcitestyle{authoryear,round}
  \newcommand{\imprimirbibliografia}{%
    \renewcommand{\bibname}{Referencias bibliográficas}%
    \bibliography{bibliografia/referencias}}
\fi

% ---- Hipervínculos (siempre al final del preámbulo) -------------
\usepackage[hidelinks,breaklinks=true]{hyperref}
\usepackage{url}
\urlstyle{same}


\begin{document}

% ---- Portada -----------------------------------------------------
% ---------------------------------------------------------------
%  Portada institucional UNIR
% ---------------------------------------------------------------
\begin{titlepage}
\thispagestyle{empty}
\begin{singlespace}
\begin{center}

\includegraphics[width=0.55\textwidth]{logo-unir.png}

\vspace{1.5cm}

{\large Universidad Internacional de La Rioja}\\[0.3cm]
{\large Escuela Superior de Ingeniería y Tecnología}\\[0.3cm]
{\large Maestría en Inteligencia Artificial}

\vspace{2.5cm}

{\LARGE\bfseries\color{unirazul} \tituloproyecto \par}

\vspace{3cm}

\begin{tabularx}{\textwidth}{@{}p{6cm}X@{}}
\toprule
Trabajo de innovación presentado por: & Francisco Javier Camacho Dorado \\
                                      & Mario Alberto Diaz Zetina \\
                                      & Ángel Zait Hernández López \\
                                      & Guillermo Sais de la Torre \\
\midrule
Fecha: & \fechaproyecto \\
\bottomrule
\end{tabularx}

\end{center}
\end{singlespace}
\end{titlepage}


% ---- Índices -----------------------------------------------------
\pagenumbering{roman}
\begin{singlespace}
\renewcommand{\contentsname}{Índice de contenidos}
\tableofcontents
\renewcommand{\listfigurename}{Índice de figuras}
\listoffigures
\renewcommand{\listtablename}{Índice de tablas}
\listoftables
\end{singlespace}

% ---- Cuerpo ------------------------------------------------------
% El salto de página lo introduce \sectionbreak (ver preambulo.tex): cada
% capítulo y cada anexo abren en hoja nueva, así que aquí no hacen falta
% \newpage adicionales.
\clearpage
\pagenumbering{arabic}

\section{Introducción}\label{introducciuxf3n}

La gestión eficiente de inventarios y la planificación logística constituyen factores fundamentales para garantizar la continuidad operativa y la calidad del servicio en organizaciones que dependen del abastecimiento periódico de productos, es el caso de aquellas que operan redes de máquinas expendedoras, donde la disponibilidad oportuna de insumos impacta directamente en la satisfacción de los usuarios, la eficiencia de los procesos de distribución y el aprovechamiento de los recursos operativos \citep{grzybowska2020,mehmood2023}.

Sin embargo, muchas organizaciones continúan gestionando sus actividades de reabastecimiento mediante esquemas basados en calendarios fijos, inspecciones periódicas o intervenciones reactivas derivadas de incidencias reportadas por los clientes \citep{rosello2023}. Este tipo de estrategias presenta diversas limitaciones, cuando el reabastecimiento se realiza siguiendo rutas predeterminadas, sin considerar el comportamiento real del consumo, se pueden generar desplazamientos innecesarios hacia equipos que aún cuentan con niveles adecuados de inventario; mientras que algunas máquinas pueden agotar sus insumos antes de la siguiente visita, ocasionando interrupciones del servicio y afectando la experiencia del cliente \citep{aguas2025}. Estos eventos se traducen en un uso ineficiente de recursos logísticos, mayores costos operativos y una disminución de la calidad del servicio, desde la percepción de los clientes \citep{grzybowska2020}.

En los últimos años, la Inteligencia Artificial (IA) y el aprendizaje automático han demostrado su capacidad para apoyar la resolución de problemas relacionados a la gestión de inventarios, el pronóstico de demanda y la optimización logística \citep{puma2025}. El análisis de datos históricos permite identificar patrones de comportamiento, estimar necesidades futuras de abastecimiento y generar información útil para la toma de decisiones \citep{mehmood2023}. Asimismo, la combinación de modelos predictivos con algoritmos de optimización posibilita la planificación de rutas más eficientes, alineadas con las necesidades reales de operación \citep{grzybowska2020}.

En Bolivia el servicio de máquinas expendedoras de alimentos y bebidas es un modelo de negocio en crecimiento y todavía la oferta es limitada, se puede encontrar este servicio en aeropuertos, terminales de buses, hospitales, oficinas, edificios corporativos y espacios públicos de las principales ciudades del país: La Paz, Cochabamba y Santa Cruz \citep{datec2024}. Son pocas empresas que ofrecen este servicio como Amarket \citep{farmacorp2024}, muchas de ellas operan mediante la modalidad de comodato o venta directa para empresas e instituciones, es el caso de Smart Vending que además se encarga de toda la logística de las máquinas \citep{smartvending2022}, otras ofrecen estos servicios como una forma de promocionar sus equipos, pues en realidad son proveedores de estos equipos, como lo hacen Vending Bolivia y Nestlé Allegra \citep{vendingbolivia2026,nestle2026}.

En Sucre, la capital constitucional de Bolivia, este servicio es ofrecido por Cafetalino, una empresa que nace el año 2018, la cual ofrece autoservicio de bebidas calientes en base a café y que se constituye en pionera de este modelo de negocio en la ciudad. Inició sus actividades con unas pocas máquinas ubicadas en el aeropuerto de Alcantarí, hospitales, oficinas bancarias y centros educativos, desde entonces ha crecido permanentemente y a la fecha cuenta con una amplia red de máquinas expendedoras ubicadas en espacios públicos cerrados, que además de los ya mencionados incluyen mercados y centros comerciales, en los cuales por sus características, se tiene alguna restricción de horario; sin embargo, cuenta también con equipos ubicados en espacios abiertos, directamente al paso, hacia la acera, de manera que están disponibles las 24 horas del día.

En este contexto se desarrolla el presente proyecto, enfocado en la gestión de la red de máquinas dispensadoras de Cafetalino, empresa que de acuerdo con la información obtenida mediante entrevistas realizadas al gerente de la organización (Anexo~\ref{anexo-entrevista-al-gerente-de-cafetalino}) y a un operador responsable de las actividades de reabastecimiento (Anexo~\ref{anexo-entrevista-a-un-operador-de-cafetalino}), enfrenta situaciones recurrentes asociadas a visitas de reposición que resultan innecesarias, modificaciones extraordinarias de rutas y episodios de desabastecimiento que afectan la continuidad del servicio, condiciones que generan costos adicionales de operación y repercuten negativamente en la experiencia de los usuarios finales, debido a que maneja un sistema de reabastecimiento periódico programado de insumos y en ocasiones reactivo al recibir reportes de fallos en el servicio; además de no disponer de un indicador que le permita saber si una determinada máquina, necesita o no recarga de insumos.

Por lo expuesto, el problema central se puede enunciar de la siguiente manera:

El actual sistema de reabastecimiento periódico y de rutas fijas utilizado por Cafetalino genera un impacto negativo doble: por un lado, máquinas desabastecidas que prestan un servicio inadecuado y generan sentimientos negativos en sus clientes; y por otro, un gasto innecesario de tiempo y recursos logísticos al realizar recorridos fijos de revisión a equipos que no siempre requieren reposición de insumos.

Para atender esta problemática, se propone el desarrollo de un sistema inteligente de apoyo a la toma de decisiones basado en técnicas de aprendizaje automático y optimización logística. Esta solución utilizará los registros históricos de operación disponibles en Cafetalino para estimar las necesidades futuras de reposición de insumos de cada máquina y realizar la planificación dinámica de recorridos de abastecimiento; con ello, se busca mejorar la eficiencia operativa, reducir desplazamientos innecesarios y contribuir a la disminución de incidencias relacionadas con el desabastecimiento.

El alcance se limita al desarrollo de un Producto Mínimo Viable (MVP) enfocado en la predicción de los requerimientos de reabastecimiento y en la planificación de las rutas de recarga. La propuesta no contempla sensores físicos, dispositivos IoT ni monitoreo en tiempo real, sino que aprovecha la información histórica existente como fuente principal de conocimiento; esta delimitación responde a las necesidades identificadas en la organización y a las restricciones técnicas y económicas evidenciadas durante el análisis.

En consecuencia, el objetivo general del trabajo es mejorar la gestión del reabastecimiento de la red de máquinas dispensadoras de Cafetalino mediante un sistema basado en Inteligencia Artificial que prediga la demanda y permita la planificación dinámica de las rutas logísticas. La contribución esperada es un efecto observable sobre la operación: menos visitas innecesarias, menos episodios de desabastecimiento y una asignación más eficiente de los recursos logísticos. El detalle de este objetivo y de los objetivos específicos en que se despliega se presenta en el Capítulo 2.

\subsection{Justificación del Trabajo}\label{justificaciuxf3n-del-trabajo}

La relevancia de la propuesta radica en que aborda una problemática real mediante el aprovechamiento de datos de operación disponibles desde el 2018, transformando dicha información en un activo estratégico para apoyar la planificación de las actividades de abastecimiento. En el contexto organizacional, se espera contribuir a una gestión más eficiente de los recursos logísticos, proporcionando información que facilite la toma de decisiones de la empresa y del personal operativo.

Además, la implementación de este sistema tendrá un impacto directo en la modernización de la logística del mercado en Sucre y en Bolivia, un entorno donde no existen soluciones de gestión predictiva para redes de distribución y que se limita a métodos reactivos y de rutas fijas: promoverá la eficiencia en el transporte, evitará traslados redundantes y reducirá los costos de distribución. Desde una perspectiva tecnológica, el proyecto permitirá evaluar la viabilidad de integrar aprendizaje automático y optimización de rutas en un entorno de operación real, mostrando cómo la inteligencia artificial puede abordar problemas complejos de la cadena de suministro adaptándose a las particularidades geográficas locales y a la volatilidad del consumo sin depender de telemetría en tiempo real.

Adicionalmente, el proyecto posee relevancia académica al constituir un caso aplicado de utilización de técnicas de Inteligencia Artificial para resolver un problema de logística y abastecimiento en un contexto empresarial real. La investigación integra conceptos de análisis de datos, aprendizaje automático y optimización operativa dentro de una propuesta orientada a generar valor práctico y conocimiento aplicable en futuros proyectos de innovación tecnológica. Asimismo, los clientes se beneficiarán con una mayor disponibilidad de productos y una reducción de incidencias asociadas a la falta de insumos.

El proyecto tiene tres beneficiarios. Los clientes, que recibirán un servicio ininterrumpido y de mejor calidad, sin los episodios de desabastecimiento que hoy generan frustración y desconfianza. Los operadores de reabastecimiento, que dejarán de realizar recorridos innecesarios o redundantes y aprovecharán mejor su jornada. Y Cafetalino, que convertirá su registro histórico de recargas en un activo de alto valor y reducirá sus costos de operación al eliminar los gastos logísticos superfluos.

\subsection{Estructura en secciones}\label{estructura-en-secciones}

El presente documento se estructura de la siguiente manera: el Capítulo 1 presenta el contexto, la problemática y la justificación del proyecto; el Capítulo 2 establece el objetivo general y los objetivos específicos; el Capítulo 3 desarrolla la aplicación de Design Thinking, desde la investigación con los involucrados hasta la selección del prototipo; el Capítulo 4 describe la metodología de trabajo basada en Scrum y la filosofía LEAN; el Capítulo 5 detalla la implementación de la propuesta, con su planificación, presupuesto, construcción, despliegue y mantenimiento; el Capítulo 6 plantea el diseño experimental con el que se validará la solución en operación; y el Capítulo 7 contrasta los resultados con los objetivos planteados y señala las líneas de trabajo futuro. Los anexos recogen los instrumentos de investigación, los diagramas de los prototipos, las capturas de la gestión del proyecto y el backlog inicial.

\section{Objetivos}\label{objetivos}

Este apartado constituye el puente entre el problema identificado en el capítulo anterior y la contribución que se propone. El esquema de reabastecimiento periódico y de rutas fijas que utiliza Cafetalino provoca visitas innecesarias a equipos que aún disponen de insumos y, al mismo tiempo, episodios de desabastecimiento que interrumpen el servicio, de modo que la propuesta debe traducirse en un efecto medible sobre esa operación. En consecuencia, a continuación se presentan tanto el objetivo general, como los objetivos específicos que contribuyen a su logro, planteados para el presente proyecto de innovación.

\subsection{Objetivo General}\label{objetivo-general}

Mejorar la gestión del reabastecimiento de la red de máquinas dispensadoras de Cafetalino, reduciendo las visitas innecesarias y los episodios de desabastecimiento, mediante un sistema basado en Inteligencia Artificial que prediga la demanda y permita la planificación dinámica y la optimización de las rutas logísticas, con el fin de contribuir a la continuidad del servicio, la reducción de ineficiencias operativas, mejorar la experiencia del cliente y minimizar los costos logísticos.

\subsection{Objetivos Específicos}\label{objetivos-especuxedficos}

Para el logro del objetivo general, se formularon los siguientes objetivos específicos:

\begin{itemize}
\item
  Caracterizar los patrones históricos de consumo y reabastecimiento de la red de máquinas dispensadoras de Cafetalino a partir de la información operativa disponible entre 2018 y 2026.
\item
  Desarrollar un módulo para generar estimaciones de los requerimientos futuros de reposición de insumos mediante la aplicación de técnicas de aprendizaje automático basadas en datos históricos de operación, sin la instalación de hardware adicional.
\item
  Desarrollar un módulo para optimizar la planificación de recorridos de abastecimiento mediante la utilización de rutas dinámicas, construidas a partir de la demanda proyectada para cada máquina dispensadora.
\item
  Desarrollar una interfaz web para el personal operativo de Cafetalino, que se constituya en una herramienta de apoyo a la toma de decisiones, y permita visualizar pronósticos de la demanda y las rutas de abastecimiento recomendadas.
\item
  Evaluar el desempeño de la solución propuesta mediante métricas de error de predicción y de reducción de recorridos frente al esquema de rutas fijas actual.
\end{itemize}

\section{Desarrollo conceptual}\label{desarrollo-conceptual}

El desarrollo conceptual para este proyecto de innovación se desarrolló aplicando sistemáticamente la metodología Design Thinking, que permite enfocar la innovación tecnológica en las necesidades de los usuarios: los clientes y trabajadores de Cafetalino, y la empresa como tal.

\subsection{Empatizar}\label{empatizar}

Esta fase está orientada a comprender las necesidades, dificultades y expectativas de todos los involucrados en el proceso de abastecimiento de la red de máquinas dispensadoras de Cafetalino. De acuerdo con los principios de Design Thinking, esta etapa busca construir una comprensión profunda del contexto operativo y de los problemas que experimentan los usuarios antes de plantear posibles soluciones tecnológicas.

Para la recopilación de información se empleó una estrategia de investigación cualitativa basada en entrevistas semiestructuradas y encuestas. Para las entrevistas se utilizaron guías previamente definidas que se incluyen en los anexos del presente documento y se seleccionaron dos participantes clave mediante un criterio de relevancia operativa y conocimiento del proceso analizado: el gerente propietario de Cafetalino, responsable de la toma de decisiones estratégicas y de la supervisión general del negocio (Anexo A), y el operador de reabastecimiento con mayor experiencia dentro de la empresa, responsable directo de la ejecución de las rutas de abastecimiento (Anexo B); en tanto que las encuestas se aplicaron a los clientes de Cafetalino (Anexo C), que son la razón de ser de las operaciones de la empresa y los receptores directos de los resultados de dichas operaciones.

Tanto las entrevistas como las encuestas se realizaron entre junio y julio de 2026. Para la aplicación de las encuestas se realizó un muestreo aleatorio estratificado de las maquinas en función al tipo de ubicación: Hospitales, centros de salud, centros educativos, centros comerciales, espacios abiertos hacia la acera, y un muestreo no probabilístico consecutivo a los clientes de Cafetalino que acudieron las ubicaciones seleccionadas en el periodo definido para la encuesta; sin embargo, la participación fue voluntaria y anónima, habiéndose encuestado a un total de 254 clientes que acudieron a 15 máquinas dispensadoras (Anexo C).

Un resumen de la entrevista realizada al gerente de Cafetalino se detalla a continuación: A la consulta de problemas o procesos que le gustaría solucionar u optimizar, mencionó varias ideas, una que llamo la atención por las posibilidades de implementación con inteligencia artificial es la relacionada a un sistema de gestión de inventarios que permita controlar la cantidad de insumos en las máquinas y pronosticar sus requerimientos de recarga (Anexo A).

Indicó que su empresa posee un registro de datos de sus operaciones desde el año 2018 a la fecha, que ve con buenos ojos el uso de IA para modernizar sus operaciones, y está dispuesto a compartir esta información bajo un acuerdo de confidencialidad estricto para protegerse de la competencia. Entre sus experiencias tecnológicas previas, señaló que intentó usar sensores para monitorear la cantidad de vasos en sus máquinas, pero el proyecto piloto fracasó debido a problemas de conectividad, altos costos y limitaciones técnicas para medir otros insumos principalmente sólidos y líquidos (Anexo A).

Comentó algunos conflictos con sus clientes, que se presentaron cuando una máquina quedó sin insumos, ellos experimentaron frustración, decepción, impotencia, sensación de pérdida o robo y pérdida de confianza, que alguna vez escaló en agresiones físicas a las máquinas, insultos al personal, difamaciones en redes e incluso un caso de sabotaje intencional introduciendo objetos extraños. La empresa mitiga esto reembolsando el dinero, actualmente esto se facilitó mediante el empleo de transferencias QR, y otorgando bebidas de cortesía (Anexo A).

En relación a la recarga de insumos el personal cubre dos rutas al día, pero el proceso tiene fallas: en ocasiones visitan máquinas que no necesitan recarga, o se ven obligados a alterar las rutas y priorizar emergencias cuando una máquina se queda fuera de servicio por falta de insumos, lo que exige llevar insumos extra a sus recorridos rutinarios y a veces hacer viajes exclusivos para la recarga de un solo equipo (Anexo A).

Un resumen de la entrevista al operador de recarga más antiguo de la empresa, cuyo resumen es el siguiente: El operador recorre dos rutas diarias, una por la mañana y otra por la tarde, que incluyen de 4 a 6 máquinas cada una. Antes de salir, prepara los paquetes de insumos necesarios más dos adicionales en previsión. En cada visita, el operador apaga la máquina, la limpia, le hace mantenimiento, la recarga y la enciende para verificar que no haya fugas o ruidos extraños, estas rutas están pre-definidas por zonas; sin embargo, alguna vez el operador se vio en la necesidad de modificar su ruta o cambiar la secuencia de visitas para priorizar las máquinas que son reportadas como fuera de servicio (Anexo B).

El operador señala que el procedimiento actual no es eficiente, obligándolo muchas veces a trabajar horas extra debido a los desvíos imprevistos, además reporta que se pierde mucho tiempo al visitar máquinas cuyo consumo no es el esperado y no requieren reposición de insumos ya que el mantenimiento se dificulta con los depósitos llenos. También considera que un sistema que prediga la demanda sería sumamente útil y evitaría la pérdida de tiempo al programar rutas exclusivas para máquinas reportadas fuera de servicio, minimizando los cambios de ruta a último momento (Anexo B).

Posteriormente, las respuestas obtenidas fueron revisadas y clasificadas en categorías temáticas relacionadas con la operación logística, la gestión de inventarios, la experiencia del cliente y las oportunidades de mejora mediante Inteligencia Artificial. El análisis permitió identificar patrones recurrentes coincidentes entre ambos participantes, los cuales se consideraron hallazgos relevantes para la definición del problema y el diseño de la propuesta de innovación, estas conclusiones se resumen en la siguiente tabla.

\captionof{table}{Conclusiones de las entrevistas}\label{tab:conclusiones-de-las-entrevistas}

\begin{longtable}[]{@{}
  >{\raggedright\arraybackslash}p{(\columnwidth - 4\tabcolsep) * \real{0.3333}}
  >{\raggedright\arraybackslash}p{(\columnwidth - 4\tabcolsep) * \real{0.3333}}
  >{\raggedright\arraybackslash}p{(\columnwidth - 4\tabcolsep) * \real{0.3333}}@{}}
\toprule\noalign{}
\begin{minipage}[b]{\linewidth}\raggedright
\textbf{Hallazgo identificado}
\end{minipage} & \begin{minipage}[b]{\linewidth}\raggedright
\textbf{Evidencia obtenida}
\end{minipage} & \begin{minipage}[b]{\linewidth}\raggedright
\textbf{Impacto observado}
\end{minipage} \\
\midrule\noalign{}
\endhead
\bottomrule\noalign{}
\endlastfoot
\textbf{Existen visitas de reabastecimiento que no generan valor} & Aproximadamente una de cada diez visitas se realiza a máquinas que no requieren reposición. & Incremento de tiempo operativo y costos logísticos. \\
\textbf{Se producen modificaciones extraordinarias de ruta} & Se reportan aproximadamente diez eventos de modificación de recorridos por año. & Ineficiencia en la planificación logística. \\
\textbf{El desabastecimiento afecta la experiencia del cliente} & Se registran quejas, reclamos y pérdida de confianza cuando la máquina no puede prestar el servicio esperado. & Impacto negativo en la percepción del servicio. \\
\textbf{Existen registros históricos disponibles desde 2018} & Información proporcionada por la dirección de la empresa. & Viabilidad para desarrollar modelos predictivos. \\
\textbf{Intentos previos basados en sensores no fueron exitosos} & Problemas de conectividad y costos de implementación. & Necesidad de una solución basada en datos históricos. \\
\end{longtable}

\fuente{Fuente: Elaboración propia}

Como resultado de las encuestas (Anexo C) se observa que el 60.6 \% de los clientes consume utiliza los servicios de la empresa dos o más veces a la semana, el 28.7 \% una vez por semana, luego el 89.4 \% son clientes asiduos, el 8.3 \% usa el servicio una vez al mes y el 2.0 \% son clientes esporádicos, el 0.4 \% no respondió la pregunta.

\begin{figure}[!htbp]
\centering
\caption{Hábitos de consumo de los clientes de Cafetalino}\label{fig:habitos-de-consumo-de-los}
\includegraphics[width=0.85\textwidth]{grafico1.png}
\fuente{Fuente: Elaboración propia}
\end{figure}

En relación al grado de satisfacción, el 82.7 \% de los clientes encuestados califica el servicio como bueno o muy bueno, el 6.3 \% como regular y el 9.5 \% como malo o muy malo, el 1.6 \% no respondió esta pregunta.

\begin{figure}[!htbp]
\centering
\caption{Grado de satisfacción de los clientes de Cafetalino}\label{fig:grado-de-satisfaccion-de-los}
\includegraphics[width=0.85\textwidth]{grafico2.png}
\fuente{Fuente: Elaboración propia}
\end{figure}

Consultados sobre la ocurrencia de fallos al usar las máquinas de Cafetalino, el 90.2 \% indico que nunca se les presento ningún problema, el 2.8 \% que recibió un producto que no tenía la calidad esperada y el 6.3 \% que no recibió el servicio, el 0.8 \% no respondió la pregunta. El 82.6 \% indico que le ocurrió una sola vez, y el 17.4 \% que le ocurrió en dos ocasiones.

\begin{figure}[!htbp]
\centering
\caption{Ocurrencia de fallos en el servicio de Cafetalino}\label{fig:ocurrencia-de-fallos-en-el}
\includegraphics[width=0.85\textwidth]{grafico3.png}
\fuente{Fuente: Elaboración propia}
\end{figure}

Respecto a sus sentimientos por el percance ocurrido el 91.3 \% de los afectados indicó que se sintió decepcionado, el 78.3 \% sintió frustración, un 26.1 \% impotencia, el 4.4 \% sintió que se le había robado y que perdió su dinero, al 43.5 \% le genero desconfianza y el 56.5 \% indica que a pesar del inconveniente no perdió la confianza.

\begin{figure}[!htbp]
\centering
\caption{Sentimientos del cliente de Cafetalino ante ocurrencia de fallos del servicio}\label{fig:sentimientos-del-cliente-de-cafetalino}
\includegraphics[width=0.85\textwidth]{grafico4.png}
\fuente{Fuente: Elaboración propia}
\end{figure}

El 82.6 \% de los afectados reportó el incidente a Cafetalino, el 13.0 \% no lo reporto y el 4.4 \% no respondió a esta consulta. El 89.5 \%, de los que reportaron el suceso, lo hizo llamando al número de contacto y el 10.5 \% en forma personal, el 100 \% de ellos indica que recibió una disculpa y el reembolso de su dinero, y el 64.7 \% que recibió una compensación por el inconveniente.

\begin{figure}[!htbp]
\centering
\caption{Respuesta de Cafetalino al Cliente ante fallos del servicio}\label{fig:respuesta-de-cafetalino-al-cliente}
\includegraphics[width=0.85\textwidth]{grafico5.png}
\fuente{Fuente: Elaboración propia}
\end{figure}

El 27.3 \% de los afectados volvió a utilizar el servicio de Cafetalino el mismo día o un día después, el 9.1 \% dos días después, el 36.4 \% una semana después del evento, el 18.2 \% dos semanas después, y el 9.1 \% tardo más de 2 semanas en volver a usar los servicios de la empresa.

\subsection{Definir}\label{definir}

El objetivo de esta fase consiste en sintetizar la información obtenida durante la etapa de empatizar para transformar los hallazgos identificados en un problema claramente delimitado y susceptible de ser abordado mediante una solución basada en Inteligencia Artificial.

A partir de las entrevistas realizadas al gerente de Cafetalino y al operador responsable del reabastecimiento, así como de las encuestas realizadas a los clientes de la empresa, se identificaron tres problemas principales: la existencia de recorridos que no generan valor al visitar máquinas que no requieren reposición de insumos, la necesidad de modificar rutas por eventos de desabastecimiento no previstos y el impacto negativo que la falta de productos genera en la experiencia de los usuarios. Con base en estos hallazgos se formularon los siguientes retos de diseño:

\begin{enumerate}
\def\labelenumi{\arabic{enumi}.}
\item
  ¿Cómo se podría predecir los requerimientos futuros de reposición de insumos de cada máquina dispensadora para reducir los eventos de desabastecimiento y mejorar la continuidad del servicio?
\item
  ¿Cómo se podría utilizar la demanda proyectada para priorizar únicamente las máquinas que realmente requieren recarga y reducir desplazamientos innecesarios del personal operativo?
\item
  ¿Cómo podríamos optimizar la planificación de recorridos de abastecimiento para disminuir modificaciones extraordinarias de ruta y mejorar la eficiencia logística de la empresa?
\end{enumerate}

Para evaluar el desempeño de la propuesta se establecen los siguientes criterios cuantificables de éxito, que se presentan en la siguiente tabla:

\captionof{table}{Criterios de éxito}\label{tab:criterios-de-exito}

\begin{longtable}[]{@{}
  >{\raggedright\arraybackslash}p{(\columnwidth - 4\tabcolsep) * \real{0.3247}}
  >{\raggedright\arraybackslash}p{(\columnwidth - 4\tabcolsep) * \real{0.3506}}
  >{\raggedright\arraybackslash}p{(\columnwidth - 4\tabcolsep) * \real{0.3247}}@{}}
\toprule\noalign{}
\begin{minipage}[b]{\linewidth}\raggedright
\textbf{Indicador}
\end{minipage} & \begin{minipage}[b]{\linewidth}\raggedright
\textbf{Línea base}
\end{minipage} & \begin{minipage}[b]{\linewidth}\raggedright
\textbf{Meta del proyecto}
\end{minipage} \\
\midrule\noalign{}
\endhead
\bottomrule\noalign{}
\endlastfoot
\textbf{Visitas de reabastecimiento innecesarias.} & Aproximadamente 10\% de las visitas realizadas (1 de cada 10 visitas). & Reducirlas a un máximo de 5\%. \\
\textbf{Modificaciones extraordinarias de rutas.} & Aproximadamente 10 eventos por año. & Reducirlas al menos en 50\%. \\
\textbf{Precisión del modelo predictivo.} & Método actual basado en experiencia operativa. & Alcanzar una precisión mínima de 85\%. \\
\textbf{Incidencias relacionadas con desabastecimiento.} & Registro histórico de la empresa. & Reducirlas en al menos 30\%. \\
\end{longtable}

\fuente{Fuente: Elaboración propia}

La línea base utilizada corresponde al procedimiento actual de reabastecimiento empleado por Cafetalino, el cual se fundamenta en rutas predefinidas y decisiones operativas apoyadas principalmente en la experiencia del personal. La validación del proyecto consistirá en comparar los resultados obtenidos por el sistema propuesto contra este escenario de referencia.

\subsection{Investigación de Antecedentes}\label{investigaciuxf3n-de-antecedentes}

En esta fase se debe realizar una revisión documental para construir el fundamento teórico y determinar el estado del arte relacionados al proyecto

\subsubsection{Fundamento Teórico}\label{fundamento-teuxf3rico}

El marco conceptual que incluye el conocimiento teórico que sustenta el proyecto, se desarrolló a partir de una revisión bibliográfica relacionada al campo de estudio y se detalla a continuación.

\textbf{Aprendizaje automático (Machine learning)} \citep{bergmann2026,universidadeuropea2025,mayorga2019}.

El aprendizaje automático es un área de la inteligencia artificial (IA), que se enfoca en algoritmos capaces de aprender patrones de un conjunto de datos de entrenamiento y, posteriormente, realizar inferencias precisas sobre nuevos datos. Esta capacidad de reconocer patrones permite a los modelos de aprendizaje automático tomar decisiones o hacer predicciones sin instrucciones explícitas predefinidas. Su principal premisa es que si se optimiza el rendimiento de un modelo con un conjunto de datos de tareas semejantes a los problemas del mundo real para los que se utilizará (entrenamiento), el modelo puede predicciones precisas con datos nuevos relacionados a su caso de uso final, donde el entrenamiento es sólo el medio para lograr un fin: la generalización, traducir un buen rendimiento en los datos de entrenamiento a resultados útiles en escenarios del mundo real. En esencia, un modelo entrenado aplica los patrones aprendidos a partir de los datos de entrenamiento para inferir el resultado correcto para una tarea del mundo real, la inferencia.

El aprendizaje automático funciona mediante lógica matemática, por tanto, las características relevantes o rasgos de cada dato se deben expresar numéricamente, para que un dato, un vector de características, se pueda introducir a un algoritmo matemático que aprenda a relacionar cada entrada con una salida.

Los métodos de aprendizaje automático se clasifican, según sus objetivos de entrenamiento y el tipo de datos de entrenamiento que implican, en tres paradigmas de aprendizaje: El aprendizaje supervisado, el aprendizaje no supervisado o el aprendizaje por refuerzo.

\begin{itemize}
\item
  \textbf{El aprendizaje supervisado} entrena un modelo para predecir el resultado correcto a partir de una entrada dada y se aplica a tareas que requieren cierto grado de precisión respecto a una verdad fundamental externa, como la clasificación o la regresión.
\item
  \textbf{El aprendizaje no supervisado} entrena un modelo para descifrar patrones intrínsecos, dependencias y correlaciones en los datos no requiere una referencia externa para comparar sus resultados.
\item
  \textbf{El aprendizaje por refuerzo (RL)} entrena un modelo para evaluar su entorno y tomar la acción que le reporte una recompensa mayor. Los escenarios de RL no implican la existencia de una única verdad fundamental, pero sí la existencia de acciones buenas y malas, o en ocasiones acciones neutrales.
\end{itemize}

El entrenamiento de extremo a extremo para un modelo dado puede, y a menudo lo hace, involucrar enfoques híbridos que aprovechan más de un paradigma. Por ejemplo, el aprendizaje no supervisado se usa frecuentemente para preprocesar datos que serán usados en el aprendizaje supervisado o por refuerzo. Los modelos de lenguaje a gran escala (LLM) suelen someterse a un entrenamiento inicial (pre-entrenamiento) y luego ser ajustados ajuste, afinar su entrenamiento a través de variantes de aprendizaje supervisado, seguido de un ajuste más fino a través de técnicas de aprendizaje por refuerzo (RL) a partir de retroalimentación humana (RLHF).

\textbf{Regresión} \citep{lee2026,opit2023}

La regresión destaca como una de las técnicas de aprendizaje automático supervisado, que sirve de puente entre el pasado, el presente y el futuro, que identifica diferentes eventos del pasado y los desglosa para su análisis, a partir de este análisis, extrae conclusiones sobre el futuro y ayuda a planificar los siguientes pasos. Se basa en aprender las relaciones entre las variables de entrada y las variables de salida, los objetivos, de manera que sus modelos predicen valores continuos.

Además del rendimiento predictivo, los métodos de aprendizaje supervisado tienen dos ventajas: La explicabilidad, que permite comprender por qué el modelo realiza una determinada predicción, algo que es muy útil en ámbitos regulados; y la eficiencia, pues los algoritmos clásicos suelen requerir menos recursos y datos, lo que los hace ideales para aplicaciones a pequeña o mediana escala y para sistemas de decisión integrados.

Existen diversos tipos de regresiones, entre ellos:

\begin{itemize}
\item
  \textbf{Regresión lineal simple}, una de las técnicas más comunes en el aprendizaje automático, algoritmo que recibe su nombre por su funcionamiento: analiza en profundidad la relación entre variables independientes y dependientes, para a partir de estos resultados, realizar predicciones sobre el futuro, demostrando su utilidad en diversas aplicaciones.
\item
  \textbf{Regresión polinómica}, una técnica en la que la salida se modela como una función polinómica de la entrada, que permite obtener curvas en lugar de líneas rectas.
\item
  \textbf{Regresión múltiple o regresión lineal múltiple}, donde el resultado depende de varias variables de entrada, que permite al modelo capturar el efecto combinado de varios predictores sobre el resultado.
\item
  \textbf{Regresión de cresta o regresión L2}, una versión de regresión lineal regularizada, que desalienta los coeficientes grandes penalizando sus valores al cuadrado, lo que reduce el sobreajuste cuando las características están correlacionadas.
\item
  \textbf{Regresión Lasso o regresión L1}, una técnica que no solo penaliza los coeficientes grandes, sino que puede reducir algunos a cero, útil para seleccionar solo las características más importantes.
\item
  \textbf{Regresión de red elástica}, un método híbrido que combina la regularización L1 y L2 para equilibrar la selección de características y la reducción de coeficientes, especialmente útil cuando las características están altamente correlacionadas o cuando hay más predictores que observaciones.
\item
  \textbf{Regresión logística}, un algoritmo de clasificación que utiliza una función logística para modelar la probabilidad de un resultado binario en función de las características de entrada.
\item
  \textbf{Modelos autorregresivos (AR)}, son modelos de series temporales en los que el valor actual se predice como una combinación lineal de sus valores anteriores.
\end{itemize}

La regresión comparte una familia de potentes algoritmos de aprendizaje supervisado con la y clasificación, se pueden adaptarse a ambos contextos; sin embargo, la elección del modelo y su configuración, dependen del tipo de resultado: Continuo para regresión o categórico para clasificación. Algunos de estos algoritmos son:

\begin{itemize}
\item
  \textbf{Árboles de decisión}, son modelos que realizan predicciones mediante una serie de preguntas sencillas, donde cada pregunta divide los datos en grupos más pequeños y similares, llamados hojas, de manera que el valor promedio de cada grupo se aproxime a los valores reales, reduciendo así el error de predicción. El árbol se divide continuamente hasta que alcanza un punto de parada: Un número máximo de preguntas o la falta de ejemplos para dividir, finalmente se realizan predicciones basadas en el valor promedio del grupo final, hoja. Los árboles de regresión tienen como objetivo minimizar la suma de cuadrados residuales (RSS) dentro de cada hoja.
\item
  \textbf{Bosques aleatorios}, consisten de un conjunto de árboles de decisión, cada uno entrenado con un subconjunto de datos obtenido mediante re-muestreo (bootstrap) y con selección aleatoria de características en cada división, las predicciones finales se realizan promediando entre los árboles. Los bosques aleatorios reducen la varianza al promediar diversos modelos, cada árbol es un experto ruidoso, y su agregación da como resultado un predictor robusto, donde la selección aleatoria de características garantiza la diversidad entre los árboles.
\item
  \textbf{K vecinos muy cercanos (KNN - K-nearest neighbors)}, es un método no paramétrico basado en instancias que almacena todo el conjunto de entrenamiento; es decir que, a diferencia de la regresión lineal, no presupone ninguna distribución subyacente de los datos y no requiere la estimación de parámetros, no presupone nada sobre la función subyacente, lo que lo hace flexible pero computacionalmente costoso. Para realizar una predicción a partir de una nueva entrada, encuentra los k ejemplos de entrenamiento más cercanos, calcula el promedio de sus valores y predice basándose en la similitud local, bajo la premisa: ``eres como tus vecinos más cercanas''.
\item
  \textbf{Máquinas de vectores de soporte (SVM - Support vector machines)}, son modelos que encuentran el hiperplano, un límite de decisión que ajusta un margen alrededor de los datos, con margen máximo o desviación mínima, se enfocan en los puntos de datos más difíciles de clasificar, aquellos cercanos al límite de decisión, ignoran los puntos más sencillos e intentan encontrar el separador más robusto.
\end{itemize}

\textbf{Planificación automática} \citep{navarro2026,zafra2021}

La planificación automática calcula un plan o secuencia de acciones que se debe aplicar al estado de un problema para alcanzar otro estado que satisfaga unas determinadas condiciones. Un planificador necesita tres entradas: Una descripción del estado actual, una descripción de las acciones que se pueden realizar, para cada acción, se deben indicar qué condiciones deben cumplirse para ejecutarla y sus efectos, y una descripción del objetivo que se quiere alcanzar. Se dispone de planificadores capaces de modelar problemas que requieren un alto nivel de expresividad e incluyen parámetros de control en las acciones, de este modo, puede decidir ciertos valores numéricos para cada acción.

\subsubsection{Estado del Arte}\label{estado-del-arte}

Para recopilar los trabajos previos realizados en el campo de estudio del proyecto y que permiten conocer los avances y sus limitaciones, se realizó una revisión documental, la misma que permitió plantear la innovación de esta propuesta. Un resumen de estos avances y sus alcances se detalla a continuación.

Varios autores, que han realizado diversas experiencias y desarrollado diferentes modelos de aplicación, destacan el uso de la telemetría y la inteligencia artificial para optimizar la gestión de inventarios en máquinas expendedoras de alimentos y bebidas, evidenciando su precisión y sus ventajas que incluyen: la reducción de costos, la mejora de las experiencias de los usuarios de sus equipos y la disponibilidad de información \citep{rosello2023,ping2024,vendingsoftware2026}.

José Puma y sus colaboradores analizaron como la gestión, de la cadena de suministro, se ha transformado con la integración de las redes neuronales y la inteligencia artificial (IA), para optimizar operaciones logísticas mediante análisis de grandes volúmenes de datos en tiempo real, herramientas que permiten crear sistemas que se adaptan, aprenden de los errores, responden eficazmente a fluctuaciones en la demanda y optimizan el flujo logístico mediante simulaciones avanzadas \citep{puma2025}.

Hanna Grzybowska y colaboradores destacan que la optimización de los sistemas logísticos de las máquinas expendedoras es sumamente compleja, pues incluye varios aspectos: La asignación de productos, los puntos de reabastecimiento, los límites de productos y las rutas de reabastecimiento; por lo que para abordar todas las facetas del problema, se requerirían varias técnicas como la previsión, el aprendizaje automático, la minería de datos, la optimización combinatoria y el enrutamiento de vehículos, entre otras, que se habían abordado de manera aislada \citep{grzybowska2020}.

Su propuesta se basa en un modelo de optimización-simulación, aplicado a la solución de problemas estocásticos, aunque no a los sistemas logísticos de máquinas expendedoras, llegando a la conclusión de que ofrece claros beneficios para la gestión de operaciones en este ámbito. Los resultados del caso de estudio mostraron que los ingresos netos de la red mejoraron en aproximadamente el 3,4 \%, la eficacia variaba con la popularidad de la máquina y comprobaron que el método de optimización tenía mucha más eficacia con máquinas de mayor rendimiento, con las cuales lograron mejoras de hasta el 6 \%. \citep{grzybowska2020}.

HOSTELVENDING reporta que Sensify, una startup argentina, implementó una plataforma tecnológica para optimizar la gestión de activos refrigerados en el sector de alimentos y bebidas: máquinas expendedoras, heladeras, entre otros. Su sistema recopila y analiza en tiempo real miles de datos de los puntos de venta, combinando sensores IoT, conectividad 24/7 y software de análisis predictivo para detectar patrones de consumo, optimizar rutas de mantenimiento y ofrecer alertas tempranas que evitan pérdidas de producto; además, la IA permite predecir la demanda por ubicación y momento del día, de manera que ofrece, a sus marcas asociadas, información en relación al consumo, temperatura, rotación de stock y estado de los equipos, a fin de reducir costos, evitar desperdicios y mejorar la eficiencia logística \citep{hostelvending2025}.

HOSTELVENDING menciona también que Azkoyen, empresa navarra fabricante de maquinaria, incorporó la Inteligencia Artificial (IA) en el mantenimiento de sus máquinas de café y sistemas de pago con la creación de una tecnología de aprendizaje continuo de los procesos y eventos en las máquinas, así como la interacción con las personas, para desarrollar modelos predictivos que estiman variables y generan modelos más precisos y eficientes para optimizar la calidad del servicio en las máquinas, cuyo objetivo final es mejorar la experiencia de sus usuarios en el manejo cotidiano de las mismas \citep{hostelvending2024}.

Aguas y Echeverria analizaron las causas de los quiebres de stock, a partir de un enfoque analítico para identificar patrones y predecir posibles desabastecimientos, aplicando un modelo predictivo supervisado basado en redes neuronales y el algoritmo Random Forest para anticipar las necesidades de inventario y mejorar los procesos de reposición de productos \citep{aguas2025}.

Umair Mehmood y colaboradores implementaron diversos algoritmos de aprendizaje automático, para predecir la demanda futura a partir de datos históricos de ventas de productos en máquinas expendedoras, concluyendo que el modelo XGBoost resultó ser el modelo con mejor rendimiento, además el estudio demostró la eficacia de utilizar técnicas de aprendizaje automático y características adicionales para la predicción precisa de la demanda futura de productos, para máquinas expendedoras, permitiendo una planificación estratégica del inventario, minimizar el desperdicio de existencias, mejorar la eficiencia operativa y aumentar los ingresos para la industria de las máquinas expendedoras \citep{mehmood2023}.

Como síntesis de la investigación de antecedentes realizada podemos indicar que si bien se han propuesto soluciones que utilizan inteligencia artificial para optimizar la gestión de inventarios y algunas mencionan que se puede aplicar a la optimización de rutas, la mayoría de las soluciones enfrenta estos dos aspectos de forma separada; además, las mismas dependen de sensores que mandan datos en tiempo real para predecir las necesidades de abastecimiento. En este último aspecto, es importante tomar en cuenta que Cafetalino ha explorado alternativas de solución basadas en la instalación de sensores que puedan transmitir datos en tiempo real; sin embargo, ha tropezado con limitaciones técnicas y económicas, especialmente con los insumos sólidos y líquidos, que no han permitido su implementación.

\captionof{table}{Matriz comparativa de trabajos previos realizados y la propuesta.}\label{tab:matriz-comparativa-de-trabajos-previos}

{\small\setlength{\tabcolsep}{4pt}
\begin{longtable}[]{@{}
  >{\raggedright\arraybackslash}p{(\columnwidth - 8\tabcolsep) * \real{0.1919}}
  >{\raggedright\arraybackslash}p{(\columnwidth - 8\tabcolsep) * \real{0.2366}}
  >{\raggedright\arraybackslash}p{(\columnwidth - 8\tabcolsep) * \real{0.2133}}
  >{\raggedright\arraybackslash}p{(\columnwidth - 8\tabcolsep) * \real{0.1790}}
  >{\raggedright\arraybackslash}p{(\columnwidth - 8\tabcolsep) * \real{0.1792}}@{}}
\toprule\noalign{}
\begin{minipage}[b]{\linewidth}\raggedright
\textbf{Trabajo Previo / Propuesta}
\end{minipage} & \begin{minipage}[b]{\linewidth}\raggedright
\textbf{¿Qué hace?}
\end{minipage} & \begin{minipage}[b]{\linewidth}\raggedright
\textbf{Fortalezas}
\end{minipage} & \begin{minipage}[b]{\linewidth}\raggedright
\textbf{Debilidades}
\end{minipage} & \begin{minipage}[b]{\linewidth}\raggedright
\textbf{Oportunidades de mejora}
\end{minipage} \\
\midrule\noalign{}
\endhead
\bottomrule\noalign{}
\endlastfoot
\textbf{\citep{puma2025}} & Integra IA y redes neuronales para optimizar la cadena de suministro.

Analiza grandes volúmenes de datos en tiempo real.

Realiza simulaciones avanzadas. & El sistema se adapta, aprende de los errores, responde eficazmente a fluctuaciones de demanda y optimiza el flujo logístico. & Requiere análisis de datos en tiempo real, que implica dependencia de conectividad constante. & Adaptar a entornos asíncronos donde la conectividad en tiempo real no es viable. \\
\textbf{\citep{grzybowska2020}} & Usa un modelo de optimización-simulación.

Resuelve problemas estocásticos de asignación de productos y rutas de reabastecimiento. & Aborda problemas complejos, integra previsión, ML, minería de datos y optimización combinatoria & El modelo no se aplicó a logística de máquinas expendedoras. & Aplicar este modelo específicamente a las máquinas expendedoras. \\
\textbf{Sensify \citep{hostelvending2025}} & Recopila datos en puntos de venta.

Combina sensores IoT, conectividad 24/7 y software predictivo para activos refrigerados. & Predice la demanda por ubicación y momento.

Ofrece alertas tempranas para evitar pérdidas de producto.

Optimiza rutas de mantenimiento. & Depende de sensores físicos (IoT).

Requiere conectividad ininterrumpida (24/7). & Desarrollar modelos independientes de hardware costoso y/o de conectividad física. \\
\textbf{Azkoyen \citep{hostelvending2024}} & Incorpora IA para el aprendizaje continuo de procesos y eventos en máquinas de café.

Genera modelos predictivos de mantenimiento. & Optimiza la calidad del servicio en las máquinas.

Mejora la experiencia de los usuarios en su manejo cotidiano. & Enfoque centrado en el mantenimiento y eventos de la maquinaria, en lugar de la cadena de suministro. & Vincular los datos de calidad y eventos de las máquinas con un sistema logístico integral de recarga de insumos. \\
\textbf{\citep{aguas2025}} & Aplica un modelo predictivo supervisado basado en redes neuronales y Random Forest para analizar causas de quiebres de stock. & Anticipa carencias de inventario.

Identifica patrones.

Mejora procesos de reposición de productos. & Enfocado sólo al análisis de stock.

Aborda el inventario separado del enrutamiento de vehículos. & Integrar la predicción de quiebres de stock con la planificación y optimización de rutas. \\
\textbf{\citep{mehmood2023}} & Implementa algoritmos de ML, destacando XGBoost para predecir la demanda futura.

Usando datos históricos de ventas. & Predicciones precisas que permiten una planificación estratégica, minimizan el desperdicio y mejoran la eficiencia operativa. & Limitado a predecir la demanda de productos.

No se acopla al despacho logístico. & Conectar las predicciones de demanda basadas en datos históricos con algoritmo de enrutamiento dinámico. \\
\textbf{Sistema de reabastecimiento predictivo y enrutamiento dinámico para Cafetalino} & Desarrollar un sistema predictivo bidireccional para procesar historial de recargas (2018-2026) y predecir el desabastecimiento y planificar rutas dinámicas óptimas. & No requiere hardware o telemetría adicional.

Integra predicción de la demanda y planificación de rutas en un ecosistema automatizado. & Depende sólo de datos históricos asíncronos.

No usa sensores en tiempo real, debido a limitaciones previas de conectividad y costos. & Entrenamiento continuo para mejorar el rendimiento. \\
\end{longtable}}

\fuente{Fuente: Elaboración propia}

En términos generales, el estado del arte muestra que la Inteligencia Artificial y el aprendizaje automático, mediante modelos como Redes Neuronales, Random Forest, XGBoost y algoritmos genéticos, se utilizan con éxito para predecir la demanda y optimizar la cadena de suministro en máquinas expendedoras, además expone cómo diversas investigaciones y startups comerciales han implementado sistemas que analizan patrones de consumo para anticipar el desabastecimiento de insumos, mejorar el mantenimiento de los equipos y elevar la experiencia del usuario final.

Sin embargo, a pesar de su eficacia, estas soluciones presentan dos grandes debilidades: En primer lugar, suelen abordar la gestión predictiva de inventarios y la planificación de rutas de transporte como problemas totalmente aislados, en segundo lugar, la mayoría depende de la instalación de telemetría, sensores físicos (IoT) y conectividad 24/7 para enviar datos en tiempo real, dependencia que se constituye en un obstáculo para las soluciones demandadas por Cafetalino, empresa que ya exploró este tipo de soluciones sin éxito.

En este contexto, la brecha identificada y que constituye la novedad del sistema propuesto radica en que la predicción de los requerimientos de la red de máquinas expendedoras se fundamenta en el procesamiento de datos históricos asíncronos, el registro de recargas de Cafetalino, sin la instalación de hardware, telemetría o sensores adicionales en los equipos de expendio de bebidas y encara de manera integral la planificación dinámica de rutas de reabastecimiento, eliminando la incertidumbre operativa, los problemas de conectividad y reduciendo los gastos logísticos innecesarios.

\subsection{Idear}\label{idear}

La finalidad de esta fase es generar varias opciones de solución al problema. Para ello el equipo de trabajo ha presentado las siguientes ideas de solución al problema planteado:

El equipo de trabajo ha presentado las siguientes ideas de solución al problema planteado:

\begin{enumerate}
\def\labelenumi{\arabic{enumi})}
\item
  \textbf{Asistente IA para Resolución Inmediata:} Integra un código QR en la máquina que conecta al usuario con un asistente de IA virtual capaz de gestionar reportes, pedir disculpas y emitir micro-reembolsos al instante.
\item
  \textbf{Visión Artificial de Bajo Costo:} Desarrollar una aplicación que mediante cámaras económicas dentro de los equipos tome imágenes de los contenedores, las cuales analice con algoritmos de inteligencia artificial y calcule el nivel de insumos sin necesidad de sensores complejos.
\item
  \textbf{Sistema de reabastecimiento predictivo y enrutamiento dinámico:} Entrenar un algoritmo utilizando los datos de ventas disponibles desde el 2018 a la fecha, para predecir los requerimientos de insumos de las máquinas dispensadoras de la empresa y sobre esta base planificar dinámicamente las rutas de reabastecimiento.
\item
  \textbf{Mantenimiento Predictivo por Patrones de Venta:} Desarrollar una App para monitorear el ritmo de las transacciones en tiempo real, detectar anomalías y alertar sobre posibles fallos mecánicos.
\end{enumerate}

La siguiente tabla muestra un análisis comparativo de las ideas generadas, incluyendo sus ventajas y desventajas para realizar una pre-selección de las ideas más prometedoras.

\captionof{table}{Análisis Comparativo de las Ideas Generadas}\label{tab:analisis-comparativo-de-las-ideas}

\begin{longtable}[]{@{}
  >{\raggedright\arraybackslash}p{(\columnwidth - 6\tabcolsep) * \real{0.2361}}
  >{\raggedright\arraybackslash}p{(\columnwidth - 6\tabcolsep) * \real{0.2737}}
  >{\raggedright\arraybackslash}p{(\columnwidth - 6\tabcolsep) * \real{0.2673}}
  >{\raggedright\arraybackslash}p{(\columnwidth - 6\tabcolsep) * \real{0.2228}}@{}}
\toprule\noalign{}
\begin{minipage}[b]{\linewidth}\raggedright
\textbf{Idea}
\end{minipage} & \begin{minipage}[b]{\linewidth}\raggedright
\textbf{Ventajas}
\end{minipage} & \begin{minipage}[b]{\linewidth}\raggedright
\textbf{Desventajas}
\end{minipage} & \begin{minipage}[b]{\linewidth}\raggedright
\textbf{Evaluación}
\end{minipage} \\
\midrule\noalign{}
\endhead
\bottomrule\noalign{}
\endlastfoot
\begin{minipage}[t]{\linewidth}\raggedright
\begin{enumerate}
\def\labelenumi{\arabic{enumi})}
\item
  \textbf{Asistente IA para Resolución Inmediata}
\end{enumerate}
\end{minipage} & \begin{minipage}[t]{\linewidth}\raggedright
\begin{itemize}
\item
  Mitiga sentimientos negativos del cliente.
\item
  Previene daños a las máquinas.
\item
  Automatiza reembolso a través de QR.
\end{itemize}
\end{minipage} & \begin{minipage}[t]{\linewidth}\raggedright
\begin{itemize}
\item
  Solución reactiva.
\item
  Trata el síntoma, pero no la causa.
\item
  Requiere desarrollo de pasarelas de pago.
\end{itemize}
\end{minipage} & Atiende el pilar de la experiencia de cliente y mitigación de daños a la marca, pero se aleja del alcance logístico principal.

\textbf{3/10} \\
\begin{minipage}[t]{\linewidth}\raggedright
\begin{enumerate}
\def\labelenumi{\arabic{enumi})}
\setcounter{enumi}{1}
\item
  \textbf{Visión Artificial de Bajo Costo}
\end{enumerate}
\end{minipage} & \begin{minipage}[t]{\linewidth}\raggedright
\begin{itemize}
\item
  Supera limitaciones de medir insumos sólidos y líquidos con sensores convencionales.
\item
  Ofrece datos precisos sobre el estado físico del inventario.
\end{itemize}
\end{minipage} & \begin{minipage}[t]{\linewidth}\raggedright
\begin{itemize}
\item
  Requiere instalar hardware adicional (cámaras).
\item
  Depende de conectividad constante.
\end{itemize}
\end{minipage} & Costos elevados y riesgos de conexión.

\textbf{5.5/10} \\
\begin{minipage}[t]{\linewidth}\raggedright
\begin{enumerate}
\def\labelenumi{\arabic{enumi})}
\setcounter{enumi}{2}
\item
  \textbf{Sistema de reabastecimiento predictivo y enrutamiento dinámico}
\end{enumerate}
\end{minipage} & \begin{minipage}[t]{\linewidth}\raggedright
\begin{itemize}
\item
  Aprovecha el registro histórico de datos.
\item
  No requiere instalar hardware adicional.
\item
  Supera problemas de conectividad.
\end{itemize}
\end{minipage} & \begin{minipage}[t]{\linewidth}\raggedright
\begin{itemize}
\item
  Depende de la calidad y limpieza de los datos históricos.
\item
  No reacciona en tiempo real ante eventos anómalos repentinos
\end{itemize}
\end{minipage} & La solución más factible, aprovecha un activo existente y elimina barreras tecnológicas que hicieron fracasar intentos previos.

\textbf{8/10} \\
\begin{minipage}[t]{\linewidth}\raggedright
\begin{enumerate}
\def\labelenumi{\arabic{enumi})}
\setcounter{enumi}{3}
\item
  \textbf{Mantenimiento Predictivo por Patrones}
\end{enumerate}
\end{minipage} & \begin{minipage}[t]{\linewidth}\raggedright
\begin{itemize}
\item
  Predice un posible servicio defectuoso.
\item
  Evita mala experiencia del cliente
\end{itemize}
\end{minipage} & \begin{minipage}[t]{\linewidth}\raggedright
\begin{itemize}
\item
  Requiere transmitir transacciones en tiempo real.
\item
  Requiere instalar hardware adicional.
\item
  Choca con la conectividad deficiente.
\end{itemize}
\end{minipage} & Depende de telemetría instantánea y es inviable en el diseño actual que dispone de datos históricos asíncronos.

\textbf{2/10} \\
\end{longtable}

\fuente{Fuente: Elaboración propia}

La tabla muestra que las ideas más prometedoras son la segunda y la tercera, la primera y la cuarta no atacan el problema principal, lo que lleva a descartarlas.

\subsection{Prototipado}\label{prototipado}

En esta fase se deben construir representaciones de las ideas más prometedoras para visualizarlas y entender su funcionamiento, para así evaluar cómo resuelven el problema. Por tanto, se presentan los prototipos de la segunda idea y la tercera, que de acuerdo a la comparación realizada resultaron ser las más prometedoras.

\subsubsection{Prototipo de la Segunda Idea: Visión Artificial de Bajo Costo}\label{prototipo-de-la-segunda-idea-visiuxf3n-artificial-de-bajo-costo}

La siguiente figura muestra la estructura de funcionamiento de la aplicación Visión Artificial de Bajo Costo, para la detección de los niveles de insumos en los depósitos de las máquinas dispensadoras de Cafetalino, donde se puede ver que la aplicación se ejecutará en tres capas.

La primera en cada máquina expendedora, donde una cámara toma una imagen de los contenedores de insumos del equipo, un microcontrolador gestiona el envío de la imagen, de forma asíncrona, que se transmite a través de un módulo de comunicación usando la Red Celular o WiFi

La segunda capa se ejecuta en un servidor, en la nube, donde se recibe la información y se ejecuta la inteligencia artificial, esta capa recibe la imagen enviada desde la máquina expendedora y realiza dos acciones simultaneas: Por un lado, almacena la imagen original en un repositorio de almacenamiento de Imágenes para mantener un historial visual, y por otro activa un disparador (Trigger) de análisis que envía la imagen al motor de inteligencia artificial, donde la imagen se alimenta al Servicio de Inferencia, que realiza una segmentación y estimación del nivel físico de los insumos, el modelo evalúa la imagen procesada y devuelve como resultado el \% de volumen restante en el contenedor, que guarda y actualiza automáticamente en la Base de Datos de Inventario.

La tercera capa es la interfaz con el usuario final, que toma los datos actualizados de la Base de Datos de Inventario y alimenta dos herramientas: Un panel de monitoreo o control visual (Dashboard) donde el personal puede ver el estado general y los niveles de insumos de las máquinas en tiempo real y un sistema de alertas y enrutamiento, un módulo logístico que, al detectar que el \% de volumen de una máquina baja de un umbral crítico, emite notificaciones y ayuda a programar el reabastecimiento únicamente de los equipos que lo necesitan.

\begin{figure}[p]
\centering
\caption{Estructura de la Aplicación Visión Artificial de Bajo Costo}\label{fig:estructura-de-la-aplicacion-vision}
\includegraphics[height=0.82\textheight,keepaspectratio]{image2.png}
\fuente{Fuente: Elaboración propia}
\end{figure}

\subsubsection{Prototipo de la Tercera Idea: Sistema de reabastecimiento predictivo y enrutamiento dinámico}\label{prototipo-de-la-tercera-idea-sistema-de-reabastecimiento-predictivo-y-enrutamiento-dinuxe1mico}

A continuación, se presenta la estructura del Sistema de reabastecimiento predictivo y enrutamiento dinámico, que muestra que el sistema se ejecutara en cuatro capas consecutivas.

\begin{figure}[p]
\centering
\caption{Estructura del Sistema de Reabastecimiento Predictivo y Enrutamiento Dinámico}\label{fig:estructura-del-sistema-de-reabastecimiento}
\includegraphics[height=0.82\textheight,keepaspectratio]{image3.png}
\fuente{Fuente: Elaboración propia}
\end{figure}

La primera capa, que corresponde a las fuentes de datos, almacena la información base sin procesar, que se alimenta de dos orígenes de información: El historial de Ventas, una base de datos que contiene los registros históricos de las operaciones de la empresa desde el año 2018 hasta el presente, y el Stock Actual, un repositorio con el estado del inventario de los insumos en las máquinas dispensadoras.

La segunda capa de análisis, en la cual los datos en crudo se transforman en proyecciones futuras, el sistema extrae, transforma y carga la información de las fuentes de datos, luego los datos ingresados se filtran y estructuran para eliminar inconsistencias o errores, asegurando que sean útiles para el análisis (limpieza y preparación), posteriormente el modelo de predicción de la demanda, recibe estos datos preparados, el algoritmo matemático evalúa la información, predice qué insumos se consumirán y cuándo, además de generar como resultado final la demanda proyectada.

La tercera capa de logística, utiliza las predicciones abstractas de la demanda proyectada se individualiza, calcula el nivel de inventario esperado para cada equipo específico en la red, el sistema descarta los equipos que aún tienen un nivel adecuado de insumos y selecciona únicamente aquellos que requieren reabastecimiento, generando los nodos o destinos requeridos, las paradas estrictamente necesarias, que un algoritmo de optimización de rutas, los cruza con información externa proveniente de APIs de Mapas y Tráfico y calcula el camino más corto y eficiente, emitiendo como salida las rutas sugeridas e informes de los requerimientos de reabastecimiento de cada máquina.

La cuarta capa es la interfaz de usuario, que toma los resultados de la capa anterior y los presenta a través de una aplicación web diseñada específicamente para este cometido, facilitando la ejecución del trabajo logístico sin requerir conocimientos técnicos sobre la inteligencia artificial que opera por detrás.

\subsection{Selección del Prototipo}\label{selecciuxf3n-del-prototipo}

La finalidad de esta fase es escoger la versión de prototipo que se desarrollará y validará, en función a su viabilidad técnica y a la respuesta que da al problema.

Sobre la base de la Tabla 2 que compara las ideas generadas y el análisis de los prototipos de las ideas más prometedoras, la solución más recomendable para desarrollar y validar es el Sistema de reabastecimiento predictivo y enrutamiento dinámico. Adicionalmente, se realizó el siguiente análisis:

\begin{itemize}
\item
  La segunda idea es críticamente dependiente de la instalación de hardware adicional: la cámara, el microcontrolador y el transmisor, por lo que se requiere conectividad, la misma que anteriores proyectos fue una limitante.
\item
  Los contenedores de insumos no siempre son transparentes, en general son opacos por lo que las imágenes de la cámara no serían de mucha utilidad, ya que en muchos casos no mostrarían el nivel de insumos en los contenedores.
\end{itemize}

\section{Metodología}\label{metodologuxeda}

El desarrollo del proyecto se gestionará utilizando las metodologías Scrum y LEAN. Scrum, que de hecho está basado en la filosofía LEAN, permitirá organizar y dar seguimiento a las actividades necesarias para construir el producto mínimo viable, además de permitir la mejora continua tanto del proceso de desarrollo como del equipo de trabajo, mientras que LEAN proporcionará un marco de validación orientado a verificar que la solución propuesta genera valor para la organización. La aplicación conjunta de ambas metodologías permitirá gestionar el desarrollo de manera iterativa, medir resultados mediante indicadores cuantificables y reducir el riesgo asociado a la implementación de la propuesta de innovación.

\subsection{Roles, Artefactos y Eventos}\label{roles-artefactos-y-eventos}

Scrum establece que se debe definir el rol que tendrá cada miembro del equipo, los artefactos y herramientas que se utilizarán y los eventos que permitirán revisar y hacer los cambios necesarios para garantizar la calidad del producto. Estos elementos, en el marco del proyecto, se describen a continuación.

\subsubsection{Roles}\label{roles}

La metodología define tres roles \citep{schwaber2020}: el propietario del producto (Product Owner), responsable de las decisiones que permiten lograr el máximo valor para el producto; el Scrum Master, encargado de dirigir las reuniones del equipo y de ayudar a sus miembros a mejorar su flujo de trabajo; y el equipo de desarrollo, que en cada iteración implementa los incrementos hasta completar el producto final. Tomando en cuenta que el equipo de trabajo es pequeño, formado por cuatro integrantes, se han definido los siguientes roles:

\begin{itemize}
\item
  \textbf{Product Owner}: Francisco Camacho, considerando que la propuesta beneficiará a una empresa de Bolivia y es el miembro del equipo que está más cerca y está en contacto directo y permanente con la empresa, su gerente y sus operadores.
\item
  \textbf{Scrum Master}: Mario Diaz, elegido por el equipo por sus habilidades de gestión y dirección.
\item
  \textbf{Equipo de Desarrollo}: Los cuatro miembros del equipo: Francisco Camacho, Mario Diaz, Angel Hernández y Guillermo Sais de la Torre. En el caso de Francisco y Mario, que cumplirán un doble rol, dedicarán el 75 \% de su tiempo a las tareas de desarrollo y 25 \% de su tiempo a las tareas de gestión en el rol asignado.
\end{itemize}

La Tabla~\ref{tab:roles-del-equipo-del-trabajo} resume los roles y responsabilidades del equipo de trabajo:

\captionof{table}{Roles del equipo del trabajo en el proceso de desarrollo}\label{tab:roles-del-equipo-del-trabajo}

\begin{longtable}[]{@{}
  >{\raggedright\arraybackslash}p{(\columnwidth - 4\tabcolsep) * \real{0.3268}}
  >{\raggedright\arraybackslash}p{(\columnwidth - 4\tabcolsep) * \real{0.2483}}
  >{\raggedright\arraybackslash}p{(\columnwidth - 4\tabcolsep) * \real{0.4248}}@{}}
\toprule\noalign{}
\begin{minipage}[b]{\linewidth}\raggedright
\textbf{Integrante}
\end{minipage} & \begin{minipage}[b]{\linewidth}\raggedright
\textbf{Rol}
\end{minipage} & \begin{minipage}[b]{\linewidth}\raggedright
\textbf{Responsabilidades}
\end{minipage} \\
\midrule\noalign{}
\endhead
\bottomrule\noalign{}
\endlastfoot
Francisco J. Camacho Dorado & Product Owner y Equipo de Desarrollo. & Definir prioridades; validar entregables; asegurar el alineamiento con las necesidades de Cafetalino; desarrollo de modelos analíticos; documentación técnica. \\
Mario Alberto Díaz Zetina & Scrum Master y Equipo de Desarrollo. & Coordinar las actividades del equipo; facilitar reuniones; gestionar impedimentos; investigación; análisis de datos; validación de resultados. \\
Ángel Zait Hernández López & Equipo de Desarrollo. & Desarrollo de modelos analíticos; documentación técnica. \\
Guillermo Sais de la Torre & Equipo de Desarrollo. & Investigación; análisis de datos; validación de resultados. \\
\end{longtable}

\fuente{Fuente: Elaboración propia}

\subsubsection{Artefactos}\label{artefactos}

Scrum especifica tres artefactos principales \citep{schwaber2020}: la lista de tareas pendientes del proyecto (Product Backlog), que ordena por prioridad todo lo que se requiere para construir el producto; el registro de tareas de una iteración (Sprint Backlog), que es el plan de trabajo de ese sprint; y el incremento, una versión parcial del producto que ya funciona. Los tres se describen a continuación en el marco del proyecto.

\textbf{Lista de tareas pendientes del proyecto (Product Backlog)}

La lista de tareas pendientes del proyecto se gestionará a través de Jira, como muestra la Figura \ref{fig:lista-de-tareas-pendientes-del} del Anexo~\ref{anexo-capturas-de-jira}, y reúne las historias de usuario necesarias para desarrollar el sistema propuesto para Cafetalino, cuyos componentes son: el análisis exploratorio de los datos disponibles, el modelo predictivo de reabastecimiento, el modelo de enrutamiento dinámico y la interfaz de usuario para interactuar con ambos. Cada historia se refina en elementos más pequeños y precisos, se describe, se le asigna un orden de prioridad y se estima su nivel de complejidad antes de poder incorporarse a una iteración \citep{schwaber2020}.

La lista inicial completa se recoge en la Tabla~\ref{tab:backlog-inicial-prioridad-y-nivel} del Anexo~\ref{anexo-backlog-inicial}, con las historias agrupadas por épica, sus criterios de aceptación, su prioridad y la estimación de su nivel de complejidad. Esta última se expresa en puntos de historia (Story Points, SP) siguiendo la serie de Fibonacci, un criterio que obliga a diferenciar tamaños relativos en lugar de intentar precisiones ilusorias; la asignación se realiza en consenso con todo el equipo y pondera la dificultad técnica de la tarea, la incertidumbre asociada a ella y el esfuerzo estimado. La estimación sitúa el techo de dificultad en dos historias de 8 puntos —la construcción del modelo predictivo y el cálculo de la ruta más corta—, que son las que concentran la incertidumbre técnica del proyecto; el resto se reparte entre 2 y 5 puntos, con los valores más bajos en las historias de presentación, que se apoyan en componentes ya conocidos por el equipo.

\textbf{Sprint Backlog}

El registro de tareas de cada iteración —el plan que los desarrolladores elaboran para guiar su trabajo, con el objetivo del sprint y el subconjunto de historias seleccionadas \citep{schwaber2020}— se extraerá de esa misma lista y se manejará también en Jira, tal como se aprecia en la Figura \ref{fig:registro-de-tareas-de-una}, donde las historias aparecen clasificadas según su estado: pendiente, en ejecución, en revisión y terminada. La herramienta permitirá garantizar que cada historia de usuario se ha refinado, se le ha asignado un orden de prioridad, se ha estimado su nivel de complejidad y, fundamentalmente, que cumple las definiciones de Listo del apartado \ref{definiciones-de-listocompletado-readydone}. Cada iteración se dedicará a uno de los cuatro componentes ya enunciados, en ese mismo orden: el análisis exploratorio del historial de recargas, el modelo predictivo a partir de datos históricos y sin hardware adicional, el modelo de enrutamiento dinámico construido sobre la demanda proyectada y la interfaz web que da acceso a ambos.

\textbf{Incremento}

Cada incremento es una versión parcial del producto que ya funciona, es aditivo a los anteriores y debe verificarse a fondo para garantizar que todos operen en conjunto \citep{schwaber2020}. En el sistema propuesto, cada uno de los cuatro componentes anteriores será un incremento, que deberá cumplir las definiciones de completado del apartado \ref{definiciones-de-listocompletado-readydone} y ser integrado y probado junto con los incrementos ya desarrollados y verificados.

\subsubsection{Eventos}\label{eventos}

Scrum contempla cinco eventos \citep{schwaber2020}: la iteración o Sprint, que hace de contenedor para los otros cuatro; la reunión de planificación de la iteración (Sprint Planning); la reunión diaria de Scrum (Daily Scrum), de quince minutos, para revisar avances, resolver atascos y ajustar el plan; la reunión de revisión de la iteración (Sprint Review), donde se examina el resultado y se definen los ajustes necesarios para alcanzar el producto final; y la retrospectiva (Sprint Retrospective), dedicada al desempeño del propio equipo. Todas las reuniones del proyecto se realizarán por Google Meet en horario nocturno, para no interrumpir el trabajo cotidiano de los integrantes, y se mantendrá una comunicación permanente de coordinación por WhatsApp.

Para desarrollar el sistema de reabastecimiento predictivo y enrutamiento dinámico se ha establecido que cada iteración dure dos semanas y se han proyectado cuatro iteraciones, es decir ocho semanas efectivas de trabajo entre agosto y septiembre de 2026, tal como se observa en la Figura \ref{fig:iteraciones-del-proyecto} del Anexo~\ref{anexo-capturas-de-jira}. En una iteración no se harán cambios que pongan en riesgo su objetivo o disminuyan la calidad, aunque el alcance podrá aclararse o renegociarse con el Product Owner conforme se aprenda más sobre el producto \citep{schwaber2020}.

Las cuatro reuniones serán dirigidas por Mario Diaz (Scrum Master) y atenderán los lineamientos de Francisco Camacho (Product Owner). Las historias comprometidas en la planificación se subdividirán en piezas implementables en un día como máximo, y la retrospectiva se celebrará antes de la planificación siguiente, de modo que sus acuerdos se apliquen ya en la iteración que se abre. La Tabla~\ref{tab:eventos-de-scrum-para-el} resume la frecuencia, la duración y el objetivo de cada evento.

\captionof{table}{Eventos de Scrum para el Proyecto}\label{tab:eventos-de-scrum-para-el}

\begin{longtable}[]{@{}
  >{\raggedright\arraybackslash}p{(\columnwidth - 6\tabcolsep) * \real{0.1720}}
  >{\raggedright\arraybackslash}p{(\columnwidth - 6\tabcolsep) * \real{0.2420}}
  >{\raggedright\arraybackslash}p{(\columnwidth - 6\tabcolsep) * \real{0.1401}}
  >{\raggedright\arraybackslash}p{(\columnwidth - 6\tabcolsep) * \real{0.4459}}@{}}
\toprule\noalign{}
\begin{minipage}[b]{\linewidth}\raggedright
\textbf{Evento}
\end{minipage} & \begin{minipage}[b]{\linewidth}\raggedright
\textbf{Frecuencia / Agenda}
\end{minipage} & \begin{minipage}[b]{\linewidth}\raggedright
\textbf{Duración estimada}
\end{minipage} & \begin{minipage}[b]{\linewidth}\raggedright
\textbf{Objetivo}
\end{minipage} \\
\midrule\noalign{}
\endhead
\bottomrule\noalign{}
\endlastfoot
\textbf{Sprint Planning} & Inicio de cada sprint.

Lunes 19:00 - 20:00 & 60 min & Definir las historias de usuario a desarrollar en el sprint, revisar prioridades, estimar esfuerzo y asignar responsabilidades. \\
\textbf{Daily Scrum} & Dos veces / semana.

Lunes y jueves

18:30 - 18:45 & 15 min & Revisar el avance, dificultades y coordinar el trabajo del equipo. \\
\textbf{Sprint Review} & Final de cada sprint.

Sabado 18:00 - 18:30 & 30 min & Presentar los resultados, verificar el cumplimiento de las historias de usuario comprometidas, verificar indicadores de éxito del proyecto y validar el alineamiento con los objetivos del proyecto. \\
\textbf{Sprint Retrospective} & Final de cada sprint.

Lunes 18:30 - 19:00 & 30 min & Identificar oportunidades de mejora y definir acciones correctivas para el desempeño del equipo. \\
\end{longtable}

\fuente{Fuente: Elaboración propia}

\subsection{Definiciones de Listo/Completado (Ready/Done)}\label{definiciones-de-listocompletado-readydone}

La metodología Scrum obliga a declarar explícitamente la Definición de Listo (Definition of Ready, DoR) y la Definición de Completado (Definition of Done, DoD): los criterios que establecen cuándo una historia de usuario puede empezar a desarrollarse y cuándo está terminada. Ambas son fundamentales para garantizar la calidad del software y minimizar las historias incompletas o mal implementadas, que obligan a repetir el trabajo \citep{schwaber2020,plaskonka2023}. En este apartado se formalizan las dos definiciones que se utilizarán en el proyecto.

\subsubsection{Definición de Listo (Definition of Ready - DoR)}\label{definiciuxf3n-de-listo-definition-of-ready---dor}

La Definición de Listo reúne los requisitos que una historia de usuario debe cumplir para pasar de la lista de tareas pendientes del proyecto al registro de tareas de una iteración \citep{plaskonka2023}. En el presente proyecto, una historia será considerada lista para incorporarse a una iteración cuando cumpla las condiciones que enumera la Tabla~\ref{tab:definicion-de-listo-dor-para}:

\captionof{table}{Definición de Listo (DoR) para el proyecto}\label{tab:definicion-de-listo-dor-para}

\begin{longtable}[]{@{}
  >{\raggedright\arraybackslash}p{(\columnwidth - 2\tabcolsep) * \real{0.3141}}
  >{\raggedright\arraybackslash}p{(\columnwidth - 2\tabcolsep) * \real{0.6859}}@{}}
\toprule\noalign{}
\begin{minipage}[b]{\linewidth}\raggedright
\textbf{Criterio}
\end{minipage} & \begin{minipage}[b]{\linewidth}\raggedright
\textbf{Descripción}
\end{minipage} \\
\midrule\noalign{}
\endhead
\bottomrule\noalign{}
\endlastfoot
\textbf{Objetivo claramente definido} & La funcionalidad o entregable requerido se encuentra descrito y comprendido por el equipo. \\
\textbf{Criterios de aceptación establecidos} & Existen condiciones verificables que permitan determinar el éxito de la historia de usuario. \\
\textbf{Fuente de información identificada} & Los datos, documentos o recursos necesarios se encuentran disponibles. \\
\textbf{Dependencias identificadas} & Se conocen las actividades o entregables previos necesarios para iniciar el trabajo. \\
\textbf{Esfuerzo estimado} & La historia cuenta con una estimación acordada por el equipo. \\
\textbf{Valor para el proyecto} & La actividad contribuye directamente al cumplimiento de los objetivos del MVP. \\
\end{longtable}

\fuente{Fuente: Elaboración propia}

\subsubsection{Definición de Completado (Definition of Done - DoD)}\label{definiciuxf3n-de-completado-definition-of-done---dod}

La Definición de Completado describe formalmente cuándo un incremento cumple los criterios de calidad exigidos, es decir, el conjunto de verificaciones que una historia de usuario debe superar para declararse terminada y lista para su entrega o integración \citep{schwaber2020}. En este proyecto, una historia de usuario se considerará terminada si cumple las condiciones que enumera la Tabla~\ref{tab:definicion-de-completado-dod-para}:

\captionof{table}{Definición de Completado (DoD) para el proyecto}\label{tab:definicion-de-completado-dod-para}

\begin{longtable}[]{@{}
  >{\raggedright\arraybackslash}p{(\columnwidth - 2\tabcolsep) * \real{0.2903}}
  >{\raggedright\arraybackslash}p{(\columnwidth - 2\tabcolsep) * \real{0.7097}}@{}}
\toprule\noalign{}
\begin{minipage}[b]{\linewidth}\raggedright
\textbf{Criterio}
\end{minipage} & \begin{minipage}[b]{\linewidth}\raggedright
\textbf{Descripción}
\end{minipage} \\
\midrule\noalign{}
\endhead
\bottomrule\noalign{}
\endlastfoot
\textbf{Desarrollo completado} & La funcionalidad o actividad comprometida ha sido implementada.

Una historia relacionada al modelo predictivo implica que el modelo ha sido entrenado, evaluado mediante las métricas definidas en el proyecto, documentado y validado por el equipo.

Una historia asociada a la optimización de rutas deberá generar recorridos válidos y sus resultados deberán ser consistentes con los criterios de éxito establecidos en la fase de Definir. \\
\textbf{Validación realizada} & Se verificó que los resultados cumplen los criterios de aceptación definidos. \\
\textbf{Documentación actualizada} & Los cambios realizados se encuentran documentados dentro del proyecto. \\
\textbf{Evidencia disponible} & Existen pruebas, capturas, resultados o entregables que demuestran la ejecución de la actividad. \\
\textbf{Revisión del Product Owner} & El Product Owner valida que el resultado satisface la necesidad planteada. \\
\textbf{Integración exitosa} & El entregable puede utilizarse dentro de la solución sin generar conflictos con otros componentes. \\
\end{longtable}

\fuente{Fuente: Elaboración propia}

\subsection{Principios LEAN y métricas}\label{principios-lean-y-muxe9tricas}

La metodología LEAN se empleará como complemento de Scrum para validar que el producto mínimo viable (MVP) desarrollado genere valor real para la organización. Bajo este enfoque, la construcción del MVP estará orientada a comprobar hipótesis relacionadas con la mejora de la planificación logística y la reducción de ineficiencias operativas identificadas durante la fase de análisis. El ciclo de trabajo seguirá el principio Construir--Medir--Aprender, permitiendo evaluar continuamente los resultados obtenidos y realizar ajustes cuando sea necesario.

\subsubsection{Las etapas del proceso de desarrollo de software LEAN}\label{las-etapas-del-proceso-de-desarrollo-de-software-lean}

La aplicación eficaz de LEAN al desarrollo de software comprende cinco actividades \citep{paladiy2026}: identificar el valor que el producto ofrece al cliente; mapear el flujo de valor, distinguiendo qué partes del proceso o qué funciones aportan valor y cuáles no; crear flujo, de modo que el trabajo no se acumule entre etapas ni aparezcan cuellos de botella; establecer el sistema de tracción, para que las actividades del equipo las impulse lo que el cliente valora y no lo que pudiera hacer falta en el futuro; y buscar la perfección mediante la revisión y mejora continua del propio proceso. Como Scrum se apoya en esta misma filosofía, las cinco etapas se materializan en sus eventos, tal como se describe a continuación.

\begin{itemize}
\item
  \textbf{Identificar valor.} Las historias de usuario de la lista de tareas pendientes y el orden de prioridad asignado se originan en los requerimientos de los beneficiarios del sistema —el gerente propietario de Cafetalino, el operador de recarga y los clientes—, recogidos en las entrevistas y la encuesta, y transmitidos al equipo por Francisco Camacho (Product Owner).
\item
  \textbf{Mapear el flujo.} El análisis del flujo de trabajo y los ajustes necesarios para evitar implementar funciones innecesarias se realizan en las reuniones de Scrum y en las de planificación y revisión de cada iteración, estas últimas en el contexto más amplio del producto final.
\item
  \textbf{Crear flujo.} Las reuniones de Scrum permiten detectar y resolver cuellos de botella en el día a día, pero el evento decisivo es la retrospectiva de cada iteración, que analiza el desempeño del equipo y ajusta el proceso de desarrollo.
\item
  \textbf{Establecer el sistema de tracción.} Se manifiesta en las reuniones de planificación, donde la selección y priorización de historias responden a los requerimientos del gerente propietario y del operador de recargas, canalizados a través de las directrices del Product Owner.
\item
  \textbf{Buscar la perfección.} Ocurre en las reuniones de planificación y de revisión, orientadas a la calidad del sistema, y sobre todo en las retrospectivas, donde el desempeño del equipo y el proceso de desarrollo son objeto de análisis y mejora continua.
\end{itemize}

\subsubsection{Hipótesis de valor}\label{hipuxf3tesis-de-valor}

Con base en los hallazgos obtenidos durante las fases de Empatizar y Definir, se establecieron las siguientes hipótesis de validación:

\begin{enumerate}
\def\labelenumi{\arabic{enumi}.}
\item
  El uso de modelos predictivos basados en datos históricos permitirá reducir las visitas de reabastecimiento innecesarias realizadas por el personal operativo.
\item
  La planificación de rutas basada en la demanda proyectada permitirá disminuir las modificaciones extraordinarias de recorrido ocasionadas por eventos de desabastecimiento.
\item
  La aplicación de técnicas de aprendizaje automático proporcionará una capacidad de predicción superior al método actualmente utilizado por la organización, basado principalmente en experiencia operativa y programación periódica de visitas.
\item
  La utilización del sistema propuesto contribuirá a reducir incidencias relacionadas con desabastecimiento y mejorará la continuidad del servicio ofrecido a los usuarios finales.
\end{enumerate}

\subsubsection{Métricas}\label{muxe9tricas}

En el proceso de desarrollo de software LEAN, los indicadores de rendimiento ayudan a medir el valor, el tiempo y la eficiencia de los procesos, lo que permite evitar desperdicios \citep{paladiy2026}. Para este proyecto se aplicarán las siguientes métricas:

\begin{itemize}
\item
  Tiempo promedio que se tarda en completar una historia de usuario, desde que inicia su desarrollo hasta que es aprobada por el equipo.
\item
  Porcentaje de Historias de Usuario aprobadas en la revisión de la iteración, sin requerir ajustes, ni trabajo adicional.
\item
  Porcentaje de incrementos que provocan errores de regresión; es decir cambios que rompen funcionalidades de historias de usuario completadas en anteriores iteraciones.
\item
  Para validar las hipótesis planteadas y medir la calidad del sistema desarrollado se utilizarán indicadores cuantificables alineados con los criterios de éxito de la Tabla \ref{tab:criterios-de-exito}, establecidos en la etapa Definir.
\end{itemize}

\subsubsection{Estrategia Construir -- Medir -- Aprender}\label{estrategia-construir-medir-aprender.}

En el desarrollo del sistema propuesto, la aplicación de LEAN se realizará mediante la estrategia cíclica de validación construir-medir-aprender, que es el núcleo de esta metodología y busca transformar una idea en un MVP mediante ciclos rápidos de iteración, evaluando la respuesta del usuario y decidiendo si se debe reorientar la lógica de desarrollo del proyecto o continuar con la línea de trabajo aplicada. Sus tres etapas, en el marco del proyecto, se describen a continuación.

\textbf{Construir} el producto mínimo viable, constituido por el modelo predictivo de demanda, el componente de optimización de rutas y la interfaz de visualización para apoyo a la toma de decisiones. \textbf{Medir} su desempeño mediante los indicadores de éxito de la Tabla \ref{tab:criterios-de-exito}, definidos en la etapa Definir, comparando los resultados obtenidos contra la línea base que hoy utiliza Cafetalino. Y \textbf{aprender} de esos resultados para identificar oportunidades de mejora, ajustar los parámetros de los modelos y refinar la solución antes de una posible implementación operativa.

En conjunto, Scrum aporta la estructura con la que se organizará y verificará el trabajo —roles, artefactos, eventos y las definiciones de Listo y Completado—, mientras que LEAN aporta el criterio con el que se decidirá si lo construido genera valor para Cafetalino. Ambos marcos convergen en el mismo punto: cada incremento debe poder medirse contra la operación actual de la empresa antes de darse por bueno. El siguiente capítulo describe cómo se construyen los cuatro componentes del sistema bajo esta metodología, con qué tecnologías y con qué estimación de tiempo y coste.


% Capítulos aún no redactados en el documento original:
% \section{Implementación de la propuesta}\label{implementacion-de-la-propuesta}

En el capítulo anterior se estableció el marco de trabajo con el que se organiza el desarrollo; corresponde ahora describir cómo se lleva a la práctica la propuesta dentro de la operación de Cafetalino. Este capítulo detalla la planificación temporal y económica del proyecto, la construcción de los cuatro componentes del sistema con las tecnologías empleadas y los datos que los alimentan, el plan de puesta en producción y el de mantenimiento posterior. Para sustentar la viabilidad de lo propuesto se construyó un prototipo funcional de esos componentes, de modo que las decisiones técnicas y las cifras de desempeño que siguen no son estimaciones teóricas sino resultados medidos sobre el histórico real de la empresa.

\subsection{Planificación y estimación}\label{planificacion-y-estimacion}

La planificación parte de los elementos fijados en el capítulo de metodología: una lista de tareas pendientes de veinte historias de usuario agrupadas en cinco épicas, estimada en ochenta y un puntos de historia, y cuatro iteraciones de dos semanas ejecutadas por un equipo de cuatro integrantes entre agosto y septiembre de 2026. Este apartado convierte esos elementos en un reparto concreto de trabajo, en un cronograma con hitos verificables y en un presupuesto detallado, distinguiendo el costo de construir el sistema del costo recurrente de operarlo, porque son decisiones de inversión de naturaleza distinta para la empresa.

\subsubsection{Distribución del trabajo y cronograma}\label{distribucion-del-trabajo-y-cronograma}

El reparto de historias entre las cuatro iteraciones respeta las dependencias técnicas antes que la prioridad nominal, de manera que ninguna historia se planifica antes de que exista el insumo que consume: el origen de datos precede al modelo, porque este se entrena sobre la base que aquel construye, y la interfaz precede al enrutamiento, porque es la que produce la selección de máquinas que el optimizador convierte en recorrido. El resultado, que coincide con las cinco épicas de la lista de tareas pendientes, se recoge en la Tabla~\ref{tab:distribucion-del-trabajo-y-cronograma} junto con las fechas y el hito que permite declarar cerrada cada iteración conforme a la Definición de Completado.

\captionof{table}{Distribución del trabajo, cronograma e hitos de verificación}\label{tab:distribucion-del-trabajo-y-cronograma}

{\small\setlength{\tabcolsep}{4pt}
\begin{longtable}[]{@{}
  >{\raggedright\arraybackslash}p{(\columnwidth - 8\tabcolsep) * \real{0.08}}
  >{\raggedright\arraybackslash}p{(\columnwidth - 8\tabcolsep) * \real{0.17}}
  >{\raggedright\arraybackslash}p{(\columnwidth - 8\tabcolsep) * \real{0.29}}
  >{\raggedright\arraybackslash}p{(\columnwidth - 8\tabcolsep) * \real{0.10}}
  >{\raggedright\arraybackslash}p{(\columnwidth - 8\tabcolsep) * \real{0.36}}@{}}
\toprule\noalign{}
\textbf{Iter.} & \textbf{Fechas (2026)} & \textbf{Incremento entregable} & \textbf{Hist. (SP)} & \textbf{Hito de verificación} \\
\midrule\noalign{}
\endhead
\bottomrule\noalign{}
\endlastfoot
\textbf{1} & 3 al 14 de agosto & Base de datos del histórico de recargas y análisis exploratorio. & 1 a 4 (16) & La base se reconstruye por completo desde el archivo de origen y el análisis reporta las anomalías detectadas. \\
\textbf{2} & 17 al 28 de agosto & Modelos de predicción de los cinco insumos con sus intervalos. & 5 a 8 (21) & Informe de métricas por validación cruzada y cobertura empírica de los intervalos. \\
\textbf{3} & 31 de agosto al 11 de septiembre & Interfaz web de consulta y selección de máquinas. & 9 a 14 (22) & El operador construye la lista de recarga del día sin asistencia técnica. \\
\textbf{4} & 14 al 27 de septiembre & Optimizador de recorrido, integración y documentación técnica. & 15 a 20 (22) & Se obtiene un recorrido completo desde la selección de máquinas y se entregan los manuales. \\
\end{longtable}}

\fuente{Fuente: Elaboración propia}

La Tabla~\ref{tab:asignacion-de-historias-por-sprint} detalla, historia por historia, la composición de cada iteración; la descripción completa de cada una, con sus criterios de aceptación, se encuentra en la Tabla~\ref{tab:backlog-inicial-prioridad-y-nivel} del \ref{anexo-backlog-inicial}.

\captionof{table}{Asignación de historias de usuario por sprint}\label{tab:asignacion-de-historias-por-sprint}

{\small\setlength{\tabcolsep}{4pt}
\begin{longtable}[]{@{}
  >{\raggedright\arraybackslash}p{(\columnwidth - 4\tabcolsep) * \real{0.10}}
  >{\raggedright\arraybackslash}p{(\columnwidth - 4\tabcolsep) * \real{0.78}}
  >{\raggedright\arraybackslash}p{(\columnwidth - 4\tabcolsep) * \real{0.12}}@{}}
\toprule\noalign{}
\textbf{Sprint} & \textbf{Historias asignadas} & \textbf{Total} \\
\midrule\noalign{}
\endhead
\bottomrule\noalign{}
\endlastfoot
\textbf{1} & HU1 Carga del histórico en base de datos (5), HU2 Registro de la capacidad de depósitos por máquina (3), HU3 Análisis exploratorio de los datos de consumo (3), HU4 Tolerancia del cargador a un CSV actualizado (5). & 16 SP \\
\textbf{2} & HU5 Modelo de predicción de consumo por insumo (8), HU6 Personalización con el historial propio de cada máquina (5), HU7 Evaluación por validación cruzada (3), HU8 Predicciones masivas para todas las máquinas (5). & 21 SP \\
\textbf{3} & HU9 Interfaz en español (2), HU10 Selección de fecha objetivo (5), HU11 Stock restante en cantidad y porcentaje (3), HU12 Rango de predicción bajo/alto (5), HU13 Días desde la última recarga (2), HU14 Selección de máquinas por arrastre (5). & 22 SP \\
\textbf{4} & HU15 Envío de la selección al enrutamiento (3), HU16 Cálculo de la ruta más corta (8), HU17 Tiempo de servicio en la ruta (3), HU18 Ubicación actual como punto de partida (3), HU19 Documentación de las variables evaluadas (3), HU20 Manuales de usuario y documentación técnica (2). & 22 SP \\
\end{longtable}}

\fuente{Fuente: Elaboración propia}

La carga de la primera iteración es deliberadamente menor, porque es aquella en la que el equipo se enfrenta por primera vez al dato real y afloran las anomalías del histórico; reservar holgura allí evita arrastrar retrasos al resto del calendario. Las tres siguientes se mantienen en un rango estrecho, lo que permite tratar la velocidad de la primera como línea base y ajustar el compromiso de las restantes en cada reunión de planificación. El esfuerzo total se estima a partir de la dedicación declarada de veinte horas semanales por integrante, compatible con la actividad laboral de sus miembros: sobre ocho semanas resultan seiscientas cuarenta horas de equipo, de las cuales ochenta corresponden a las tareas de gestión que asumen el Product Owner y el Scrum Master y quinientas sesenta a construcción efectiva. Esa proporción sitúa el costo unitario en unas siete horas de trabajo por punto de historia, referencia útil para dimensionar ampliaciones futuras del alcance.

Para verificar que la carga asignada es sostenible, la Tabla~\ref{tab:capacidad-por-sprint} contrasta la capacidad disponible de cada iteración de dos semanas con los puntos de historia comprometidos. Las seiscientas cuarenta horas de equipo se reparten en partes iguales entre las cuatro iteraciones, ciento sesenta horas cada una, de las cuales veinte corresponden en promedio a la gestión del Product Owner y el Scrum Master, dejando ciento cuarenta horas de construcción efectiva por sprint.

\captionof{table}{Capacidad disponible y velocidad requerida por sprint}\label{tab:capacidad-por-sprint}

{\small\setlength{\tabcolsep}{4pt}
\begin{longtable}[]{@{}
  >{\raggedright\arraybackslash}p{(\columnwidth - 8\tabcolsep) * \real{0.12}}
  >{\raggedright\arraybackslash}p{(\columnwidth - 8\tabcolsep) * \real{0.22}}
  >{\raggedright\arraybackslash}p{(\columnwidth - 8\tabcolsep) * \real{0.22}}
  >{\raggedright\arraybackslash}p{(\columnwidth - 8\tabcolsep) * \real{0.44}}@{}}
\toprule\noalign{}
\textbf{Sprint} & \textbf{Capacidad (h)} & \textbf{Carga (SP)} & \textbf{Velocidad requerida} \\
\midrule\noalign{}
\endhead
\bottomrule\noalign{}
\endlastfoot
1 & 140 & 16 & 8,75 h/SP; holgura frente al promedio, reservada para las anomalías del primer contacto con el dato real. \\
2 & 140 & 21 & 6,67 h/SP; exige ya la velocidad cercana al promedio del proyecto. \\
3 & 140 & 22 & 6,36 h/SP; sostenida gracias a los componentes reutilizados de sprints anteriores. \\
4 & 140 & 22 & 6,36 h/SP; sostenida por la misma razón, sobre la integración final. \\
\textbf{Total} & \textbf{560} & \textbf{81} & \textbf{6,91 h/SP}, coherente con el costo unitario de siete horas por punto de historia. \\
\end{longtable}}

\fuente{Fuente: Elaboración propia}

La velocidad requerida decrece de un sprint a otro porque cada iteración reutiliza los artefactos —base de datos, procesos de validación, componentes de interfaz— construidos en las anteriores, lo que reduce el esfuerzo marginal por punto de historia; esta tendencia es consistente con la holgura deliberada de la primera iteración y con el promedio de siete horas por punto de historia ya declarado para el proyecto completo.

\subsubsection{Presupuesto}\label{presupuesto}

El presupuesto distingue la inversión inicial de construcción del gasto recurrente de operación, según detalla la Tabla~\ref{tab:presupuesto-del-proyecto}. El costo de personal valora cada perfil a tarifas de mercado del sector de tecnologías de la información; debe advertirse que se trata de estimaciones propias de referencia y no de remuneraciones efectivamente pagadas, dado que el sistema se desarrolló en el marco de un trabajo académico, y se incluyen porque son el costo que la empresa asumiría para reproducir el proyecto en condiciones de mercado. Las cifras se expresan en bolivianos y en su equivalente en dólares al tipo de cambio oficial de 6,96 bolivianos, dado que los servicios de nube y cartografía se facturan en esa moneda.

\captionof{table}{Presupuesto de construcción y costo mensual de operación}\label{tab:presupuesto-del-proyecto}

{\small\setlength{\tabcolsep}{4pt}
\begin{longtable}[]{@{}
  >{\raggedright\arraybackslash}p{(\columnwidth - 8\tabcolsep) * \real{0.38}}
  >{\raggedright\arraybackslash}p{(\columnwidth - 8\tabcolsep) * \real{0.12}}
  >{\raggedright\arraybackslash}p{(\columnwidth - 8\tabcolsep) * \real{0.16}}
  >{\raggedright\arraybackslash}p{(\columnwidth - 8\tabcolsep) * \real{0.17}}
  >{\raggedright\arraybackslash}p{(\columnwidth - 8\tabcolsep) * \real{0.17}}@{}}
\toprule\noalign{}
\textbf{Concepto} & \textbf{Horas} & \textbf{Tarifa (Bs/h)} & \textbf{Subtotal (Bs)} & \textbf{Subtotal (USD)} \\
\midrule\noalign{}
\endhead
\bottomrule\noalign{}
\endlastfoot
\multicolumn{5}{@{}l}{\textbf{Inversión inicial de construcción}} \\
Científico de datos (modelado predictivo) & 160 & 90 & 14.400 & 2.069 \\
Desarrollador de backend y optimización & 160 & 80 & 12.800 & 1.839 \\
Desarrollador de frontend & 160 & 75 & 12.000 & 1.724 \\
Analista de datos y validación & 160 & 70 & 11.200 & 1.609 \\
Gestión del producto y del proceso & 80 & 85 & 6.800 & 977 \\
Servicios de nube y cartografía durante el desarrollo & & & 1.400 & 201 \\
Licencia de gestión del proyecto (4 usuarios, 2 meses) & & & 980 & 141 \\
\textbf{Total de la inversión inicial} & \textbf{720} & & \textbf{59.580} & \textbf{8.560} \\
\midrule\noalign{}
\multicolumn{5}{@{}l}{\textbf{Costo mensual de operación}} \\
Ejecución en contenedores, con escalado a cero fuera de horario \citep{cloudrun2026} & & & 174 & 25 \\
Consultas a la matriz de distancias, una por ruta planificada \citep{distancematrix2026} & & & 139 & 20 \\
Almacenamiento de imágenes, secretos y respaldos & & & 70 & 10 \\
Dominio y certificado, prorrateo anual & & & 35 & 5 \\
\textbf{Total mensual} & & & \textbf{418} & \textbf{60} \\
\end{longtable}}

\fuente{Fuente: Elaboración propia, sobre los tarifarios vigentes de los servicios citados. Tipo de cambio referencial de 6,96 Bs/USD}

El total de horas asciende a setecientas veinte y no a seiscientas cuarenta porque las de gestión se contabilizan en línea propia, separadas de las de desarrollo de los dos integrantes que asumen el doble rol. El gasto recurrente, en cambio, es reducido: se trata de un sistema de uso interno, consultado por unos pocos operadores varias veces al día, cuyo consumo se sitúa cerca de los umbrales gratuitos de los servicios empleados. Puesto en perspectiva, equivale al costo de unas pocas visitas de recarga innecesarias al mes, de modo que la meta de reducirlas a la mitad, fijada entre los criterios de éxito de la Tabla~\ref{tab:criterios-de-exito}, basta por sí sola para sostener la operación del sistema.

\subsection{Implementación}\label{implementacion}

Este apartado describe cómo se construye el sistema: la arquitectura que lo organiza, las tecnologías que lo componen, el tratamiento que reciben los datos de origen y la lógica de cada componente. La exposición sigue el recorrido del dato, desde el registro histórico de recargas hasta el itinerario que finalmente ve el operador en el mapa.

\subsubsection{Arquitectura y pila tecnológica}\label{arquitectura-y-pila-tecnologica}

La solución conserva la arquitectura en cuatro capas definida durante el prototipado, que recoge la Figura~\ref{fig:estructura-del-sistema-de-reabastecimiento}: una capa de datos de origen, una de análisis que transforma el histórico en demanda proyectada, una de decisión logística que convierte esa proyección en un recorrido y una de presentación que la pone ante el operador. La separación es funcional y también física, ya que el análisis y la decisión residen en un servicio de aplicación mientras que la presentación es una aplicación web que se comunica con aquel mediante una interfaz de programación.

\begin{figure}[!htbp]
\centering
\includegraphics[height=0.78\textheight,keepaspectratio]{image3.png}
\caption{Estructura del Sistema de Reabastecimiento Predictivo y Enrutamiento Dinámico}\label{fig:estructura-del-sistema-de-reabastecimiento}
\fuente{Fuente: Elaboración propia}
\end{figure}

Sobre esa estructura, la Tabla~\ref{tab:tecnologias-empleadas-por-capa} enumera las tecnologías elegidas para cada capa, junto con la razón que motivó cada elección.

\captionof{table}{Tecnologías empleadas en cada capa del sistema}\label{tab:tecnologias-empleadas-por-capa}

{\small\setlength{\tabcolsep}{4pt}
\begin{longtable}[]{@{}
  >{\raggedright\arraybackslash}p{(\columnwidth - 6\tabcolsep) * \real{0.19}}
  >{\raggedright\arraybackslash}p{(\columnwidth - 6\tabcolsep) * \real{0.25}}
  >{\raggedright\arraybackslash}p{(\columnwidth - 6\tabcolsep) * \real{0.28}}
  >{\raggedright\arraybackslash}p{(\columnwidth - 6\tabcolsep) * \real{0.28}}@{}}
\toprule\noalign{}
\textbf{Capa} & \textbf{Componente} & \textbf{Tecnología} & \textbf{Justificación} \\
\midrule\noalign{}
\endhead
\bottomrule\noalign{}
\endlastfoot
Datos de origen & Base de datos del histórico y guiones de carga & SQLite y pandas & Volumen de miles de registros, que no justifica un gestor servidor. \\
Análisis & Modelos de predicción de los cinco insumos & scikit-learn \citep{pedregosa2011} & Métodos de conjunto sobre árboles, adecuados a un histórico pequeño y tabular. \\
Decisión logística & Matriz de tiempos y optimizador de recorrido & Google Distance Matrix y OR-Tools \citep{ortools2026} & Tiempos de conducción reales y un resolutor de rutas maduro y de uso libre. \\
Servicio de aplicación & Interfaz de programación & Python, FastAPI y Pydantic \citep{fastapi2026} & Validación declarativa de los mensajes y documentación automática. \\
Presentación & Aplicación web de apoyo a la decisión & React, TypeScript y Vite & Interfaz reactiva con verificación de tipos en compilación. \\
\end{longtable}}

\fuente{Fuente: Elaboración propia}

La pila es deliberadamente conservadora. Todos los componentes son de uso libre y ampliamente adoptados, lo que reduce el riesgo de dependencia de un proveedor único y facilita que la empresa encuentre en el mercado local perfiles capaces de mantener el sistema. La única dependencia comercial es el servicio de cartografía, confinada a un solo módulo, de modo que sustituirla no obligaría a modificar el resto del sistema.

\subsubsection{Origen y preparación de los datos}\label{origen-y-preparacion-de-los-datos}

El insumo del sistema es el registro histórico de recargas descrito en el apartado~\ref{datos-disponibles-y-diseno-de-la-evaluacion}. En su forma original comprende 3.646 filas que abarcan de octubre de 2018 a julio de 2026, de las cuales 2.467 corresponden a máquinas activas; las restantes documentan equipos retirados de servicio y se excluyen porque su comportamiento ya no describe la red actual. Las columnas del archivo están rotuladas en español y se traducen a identificadores en inglés durante la carga, de manera que el idioma sobreviva en los valores pero no en el código.

La decisión de diseño más relevante de esta etapa es la forma de identificar una máquina. El archivo contiene treinta y seis identificadores de equipo y dieciocho denominaciones de ubicación, pero solo diecisiete pares de coordenadas. La discrepancia se explica porque los equipos rotan entre emplazamientos y porque una misma ubicación se escribe de formas distintas según quién registre la visita. Como lo que determina el consumo es el emplazamiento y no el número de serie del equipo instalado en él, las ubicaciones se consolidan agrupando por coordenada geográfica y adoptando como denominación canónica la variante más frecuente en ese punto. Resultan así diecisiete ubicaciones canónicas, cada una con su capacidad de depósito registrada: veinte litros de agua y mil seiscientos gramos de mezcla, salvo dos emplazamientos con depósitos reducidos de novecientos gramos.

De las 2.467 lecturas activas se descartan veintiocho antes del entrenamiento: diecisiete son la primera visita registrada de cada ubicación, que por definición carece de historial previo con el que construir las variables derivadas, y once presentan un intervalo desde la recarga anterior no registrado. Estas filas se eliminan en lugar de imputarse, porque imputar el historial de una máquina equivaldría a inventar el comportamiento que precisamente se quiere predecir. Quedan 2.439 observaciones de entrenamiento.

\subsubsection{Construcción del modelo predictivo}\label{construccion-del-modelo-predictivo}

Cada insumo recibe un modelo independiente, decisión tomada porque las magnitudes difieren en escala y dinámica: obligar a un único modelo a predecir simultáneamente mililitros de agua, unidades de vaso y gramos de tres mezclas lo forzaría a un compromiso que ninguno de los cinco insumos agradece. Sobre las variables de partida descritas en el capítulo anterior se construyen veinticuatro variables de entrada, que combinan la ubicación, el intervalo desde la recarga previa, el día de la semana y el mes, y para cada insumo su media histórica, su media reciente y ambas expresadas como consumo por día.

La regla que gobierna esas variables derivadas es que ninguna observación puede ver su propio resultado: todas las medias y tasas se calculan desplazando la serie una posición antes de acumular o promediar, de modo que cada fila solo incorpora información de visitas anteriores. Por el mismo motivo se descartaron dos columnas disponibles y a primera vista informativas, el porcentaje de agua restante y el número de vasos, porque ambas se miden durante la misma visita cuyo resultado se pretende anticipar: emplearlas exigiría estar ya frente a la máquina para saber qué necesita, lo que anula el propósito de predecir antes de planificar la ruta.

El modelo seleccionado es una regresión por potenciación del gradiente sobre histogramas, en la variante propuesta por \citet{friedman2001}, elegida tras compararla mediante búsqueda aleatoria de hiperparámetros con un bosque aleatorio, frente al cual resultó superior en los cinco insumos. La búsqueda explora treinta combinaciones evaluadas por validación cruzada de cinco pliegues, el mismo procedimiento con el que se reportan las métricas. Para acompañar cada estimación con una banda de incertidumbre se entrenan además dos modelos de regresión cuantílica \citep{koenker2001} por insumo, correspondientes a los percentiles quinto y nonagésimo quinto. La Tabla~\ref{tab:desempeno-de-los-modelos} presenta el desempeño obtenido.

\captionof{table}{Desempeño de los modelos por insumo, en validación cruzada de cinco pliegues}\label{tab:desempeno-de-los-modelos}

{\small\setlength{\tabcolsep}{4pt}
\begin{longtable}[]{@{}
  >{\raggedright\arraybackslash}p{(\columnwidth - 8\tabcolsep) * \real{0.28}}
  >{\raggedright\arraybackslash}p{(\columnwidth - 8\tabcolsep) * \real{0.14}}
  >{\raggedright\arraybackslash}p{(\columnwidth - 8\tabcolsep) * \real{0.20}}
  >{\raggedright\arraybackslash}p{(\columnwidth - 8\tabcolsep) * \real{0.19}}
  >{\raggedright\arraybackslash}p{(\columnwidth - 8\tabcolsep) * \real{0.19}}@{}}
\toprule\noalign{}
\textbf{Insumo} & \textbf{Unidad} & \textbf{Error absoluto medio} & \textbf{Coef. de determinación} & \textbf{Cobertura del intervalo} \\
\midrule\noalign{}
\endhead
\bottomrule\noalign{}
\endlastfoot
Agua embotellada & mililitros & 1.702,78 & 0,371 & 84,2 \% \\
Vasos & unidades & 9,46 & --- & 84,2 \% \\
Mezcla de café & gramos & 22,33 & 0,354 & 83,5 \% \\
Mezcla de chocolate & gramos & 86,40 & 0,274 & 83,8 \% \\
Mezcla de capuchino & gramos & 90,41 & 0,308 & 83,2 \% \\
\end{longtable}}

\fuente{Fuente: Elaboración propia. El coeficiente de determinación de los vasos se omite porque el valor registrado coincide exactamente con el del agua embotellada, coincidencia que apunta a un error de transcripción en el informe de métricas; en su lugar se reporta el error absoluto medio de ese insumo, verificado de forma independiente}

Procede contrastar estos resultados con la meta comprometida entre los criterios de éxito de la Tabla~\ref{tab:criterios-de-exito}, que fijaba un coeficiente de determinación igual o superior a 0,30 en cada uno de los cinco insumos. El umbral se alcanza en el agua embotellada, la mezcla de café y la de capuchino, y la mezcla de chocolate queda marginalmente por debajo, en 0,274; el insumo de los vasos no admite pronunciamiento por la razón consignada al pie de la tabla. El criterio se cumple, por tanto, de forma parcial, y conviene decirlo antes que sortearlo, porque la causa del déficit no es de implementación sino informativa y afecta por igual a los cinco insumos.

Los coeficientes de determinación, situados entre 0,27 y 0,37, indican en efecto que una parte sustancial de la variabilidad del consumo permanece sin explicar. El resultado es genuino y conviene interpretarlo correctamente: el histórico registra cuándo y cuánto se repuso, pero no la afluencia de personas, el clima ni los eventos locales que la determinan, de modo que el techo lo impone la información disponible y no el algoritmo escogido. El chocolate, además, es el insumo de rotación más irregular de los cinco, lo que explica que sea el primero en acusar esa carencia. La comparación pertinente, en todo caso, no es contra un predictor perfecto sino contra la práctica actual, basada en un calendario fijo y en la experiencia del operador; frente a esa línea base el error absoluto medio del modelo es inferior en los cinco insumos, que es la segunda mitad del criterio comprometido y la que sí se satisface por completo.

Merece registrarse la calibración de las bandas de incertidumbre. Inicialmente se emplearon los percentiles décimo y nonagésimo, que nominalmente delimitan un intervalo del ochenta por ciento, pero su cobertura empírica medida sobre pliegues no vistos resultó de apenas el setenta por ciento, una subcobertura conocida de la regresión cuantílica con conjuntos pequeños. En lugar de publicar un intervalo mal etiquetado se ampliaron deliberadamente los percentiles para compensar el sesgo medido, con lo que la cobertura efectiva se situó en torno al ochenta y tres por ciento, cerca del valor pretendido. El intervalo que ve el operador es, por tanto, verificado y no una declaración nominal. Conviene consignar también tres ensayos que no prosperaron, porque delimitan hasta dónde llega el enfoque: sustituir el día de la semana por la proporción de días laborables del intervalo no mejoró ningún insumo; añadir una variable de tendencia a corto plazo empeoró los cinco, presumiblemente por tratarse de un estimador demasiado ruidoso; y agrupar el mes en trimestres degradó el desempeño casi tanto como eliminarlo, lo que indica que la señal estacional relevante opera a granularidad fina. Ninguno de los tres cambios se incorporó.

\subsubsection{De consumo previsto a existencias restantes}\label{de-consumo-previsto-a-existencias-restantes}

Los modelos predicen consumo, pero la decisión que el operador necesita tomar se formula en términos de existencias: no le sirve saber cuántos gramos se consumirán, sino cuánto quedará en el depósito. El sistema realiza esa conversión restando el consumo previsto de la capacidad registrada para esa máquina concreta y acotando el resultado entre cero y la capacidad total, ya que ni un depósito puede quedar en negativo ni contener más de lo que admite. El valor se acompaña de su porcentaje sobre la capacidad real de esa máquina, lo que permite comparar equipos de depósito distinto sin traducir mentalmente las unidades.

La conversión exige invertir los extremos del intervalo, puesto que el escenario de menor existencia restante se corresponde con el de mayor consumo previsto. El sistema aplica esa inversión explícitamente, además de un ajuste que garantiza que el extremo inferior nunca supere a la estimación puntual ni esta al extremo superior, precaución necesaria porque los tres modelos de cada insumo se entrenan por separado y nada les impone formalmente ese orden. Los vasos son la excepción del esquema, ya que no se almacenan en un depósito de capacidad fija y se reportan como consumo previsto sin convertir.

\subsubsection{Módulo de enrutamiento dinámico}\label{modulo-de-enrutamiento-dinamico}

Seleccionadas las máquinas que requieren atención, el problema pasa a ser de orden de visita, y el módulo lo resuelve en dos etapas. Primero solicita al servicio de cartografía la matriz de tiempos de conducción entre todos los puntos implicados, lo que incorpora la red vial real en lugar de distancias en línea recta, diferencia sustancial en una ciudad de topografía irregular. Después plantea el recorrido como un problema del agente viajero de vehículo único con retorno al origen y lo resuelve con la biblioteca OR-Tools \citep{ortools2026}, planteamiento clásico dentro de la familia de problemas de enrutamiento de vehículos \citep{tothvigo2014}.

Dos elementos adaptan ese planteamiento genérico a la operación concreta de Cafetalino. El primero es que el punto de partida no es un almacén sino la ubicación actual del operador, obtenida del dispositivo y corregible manualmente en el mapa, lo que permite replanificar a mitad de jornada desde donde se encuentre. El segundo es que al tiempo de conducción se le suma un tiempo de servicio de veinticinco minutos por máquina visitada, que es lo que toma efectuar la recarga; sin ese término la duración estimada de la jornada resultaría sistemáticamente optimista y perdería utilidad para planificar. Debe señalarse con precisión el alcance de la optimización: el resolutor se configura con una heurística de construcción por arco más económico y no se le añade una fase de mejora posterior, de modo que la solución devuelta es una construcción de buena calidad y no un óptimo demostrado. Para el tamaño del problema, diecisiete máquinas como máximo más el punto de partida, la diferencia práctica es reducida y la respuesta inmediata, lo que privilegia la interacción fluida con el operador.

\subsubsection{Interfaz web de apoyo a la decisión}\label{interfaz-web-de-apoyo-a-la-decision}

La interfaz organiza la jornada del operador en tres columnas que reproducen la secuencia de su decisión: las máquinas disponibles con su pronóstico, las que ha decidido recargar y el mapa con el recorrido resultante. El operador elige una fecha objetivo, que puede ser futura, y el sistema devuelve el pronóstico de las diecisiete ubicaciones ordenadas de manera descendente por días transcurridos desde su última recarga, criterio que sitúa arriba las máquinas más rezagadas. El paso de una lista a otra se realiza arrastrando la tarjeta de la máquina, gesto que sustituye el llenado de un formulario y hace inmediata la construcción de la lista del día.

Cada tarjeta muestra las existencias restantes previstas de agua y de las tres mezclas, en unidad absoluta y en porcentaje, junto con los vasos estimados y los días desde la última recarga como indicador de la vigencia del dato. Una decisión de diseño deliberada es que la tarjeta presenta siempre el intervalo bajo y alto además de la estimación puntual: ocultar la incertidumbre produciría una falsa sensación de precisión en un modelo que, como se ha documentado, deja parte de la variabilidad sin explicar, mientras que mostrarla permite al operador reservar su criterio en los casos de banda ancha, que es el comportamiento deseable en una herramienta de apoyo a la decisión y no de sustitución de ella. Cerrada la selección, la aplicación envía las máquinas elegidas al módulo de enrutamiento y dibuja sobre el mapa el recorrido en el orden que este determina, con la duración total estimada y la lista ordenada de paradas.

\subsubsection{Requisitos de equipamiento}\label{requisitos-de-equipamiento}

Los requisitos de la solución son modestos, coherentes con el criterio, sostenido desde la selección del prototipo, de no depender de instrumentación adicional en las máquinas. El entorno de desarrollo requiere un equipo de escritorio convencional, con ocho gigabytes de memoria suficientes para el entrenamiento, que se completa en menos de diez minutos sobre el conjunto actual. El entorno de producción necesita un contenedor con uno o dos núcleos y dos gigabytes de memoria para el servicio de aplicación, más el almacenamiento de la base de datos y de los modelos entrenados, cuyo tamaño conjunto se mide en decenas de megabytes. Del lado del operador basta un teléfono con navegador y datos móviles, que es el equipamiento del que ya dispone.

Cabe precisar, para evitar una expectativa equivocada, que se evaluó explícitamente el empleo de redes neuronales y se descartó: con algo más de dos mil cuatrocientas observaciones de entrenamiento, un modelo de esa naturaleza no dispone de datos suficientes para superar a los métodos de conjunto sobre árboles y solo añadiría complejidad de infraestructura. El prototipo completo, con el código de los tres componentes y la documentación técnica del proceso de modelado, está disponible en un repositorio público en \url{https://github.com/J4rv1sX/cafetalino-project}.

\subsection{Despliegue}\label{despliegue}

La puesta en producción persigue que el sistema pueda ser operado por Cafetalino sin asistencia técnica permanente y sin exponer la información cedida por la empresa. Este apartado describe los entornos previstos, la infraestructura y las herramientas de publicación, las medidas de protección de los datos y el plan de contingencia frente a los fallos previsibles.

\subsubsection{Entornos, infraestructura y herramientas}\label{entornos-infraestructura-y-herramientas}

Se contemplan tres entornos con propósitos distintos. El de desarrollo se ejecuta en el equipo de cada integrante, con la base de datos local reconstruible desde el archivo de origen. El de preproducción replica la configuración de producción, recibe cada cambio antes que aquella y es donde se verifican los incrementos frente a la Definición de Completado. El de producción es el que utiliza el personal de Cafetalino. Los tres comparten el código y se diferencian únicamente por configuración externa —la dirección del servicio que consume la interfaz, la clave del servicio de cartografía y la lista de orígenes autorizados—, y ninguna credencial se incorpora al código ni al repositorio.

El despliegue se apoya en servicios administrados de Google Cloud, elección coherente con el uso ya existente del servicio de cartografía y que evita a la empresa la administración de servidores. El servicio de aplicación y la interfaz web se empaquetan en contenedores y se publican en un servicio de ejecución administrada que escala a cero fuera del horario de uso, característica determinante para el costo mensual estimado. Las imágenes se almacenan en un registro de artefactos y las credenciales en un gestor de secretos, del que el servicio las lee en tiempo de ejecución. La publicación se automatiza desde el repositorio: cada integración en la rama principal construye la imagen y despliega una revisión en preproducción, mientras que la promoción a producción requiere aprobación explícita del Product Owner. Como el servicio conserva las revisiones anteriores y permite dirigir el tráfico entre ellas, la publicación puede realizarse de forma progresiva y revertirse de inmediato, propiedad sobre la que se apoya buena parte del plan de contingencia.

\subsubsection{Securización y protección de los datos}\label{securizacion-y-proteccion-de-los-datos}

La información que maneja el sistema es operativa y no personal: registros de reposición de insumos por máquina y ubicación de equipos instalados en espacios públicos o comerciales. La única excepción es la posición del dispositivo del operador, empleada como punto de partida del recorrido, que se utiliza en la sesión y no se almacena. Aun así, el histórico fue cedido bajo acuerdo de confidencialidad y constituye información comercialmente sensible, ya que revela el volumen de negocio de cada emplazamiento, por lo que el repositorio permanece privado y los respaldos se conservan cifrados y con acceso restringido.

Las medidas técnicas previstas son las habituales para un sistema de esta naturaleza. El tráfico se sirve exclusivamente por canal cifrado; la interfaz de programación acepta peticiones únicamente desde los orígenes autorizados; el acceso a la aplicación exige autenticación del personal operativo, con perfiles diferenciados entre consulta y administración; y la clave del servicio de cartografía se restringe tanto por dominio de origen como por interfaz habilitada, de modo que su eventual filtración no permita un uso facturable ajeno. Los respaldos se programan con periodicidad diaria y retención mensual, suficiente dado que el sistema es reconstruible por completo a partir del archivo de origen y de la ejecución de los guiones de carga y entrenamiento.

\subsubsection{Plan de contingencias}\label{plan-de-contingencias}

El plan se concentra en los modos de fallo efectivamente observables en el sistema construido, antes que en una enumeración genérica de riesgos. La Tabla~\ref{tab:plan-de-contingencias} recoge cada uno con su forma de detección y la respuesta prevista, bajo el criterio de que el sistema degrade su servicio de manera controlada en lugar de interrumpirlo, dado que el operador debe poder ejecutar su jornada en cualquier circunstancia.

\captionof{table}{Plan de contingencias del despliegue}\label{tab:plan-de-contingencias}

{\small\setlength{\tabcolsep}{4pt}
\begin{longtable}[]{@{}
  >{\raggedright\arraybackslash}p{(\columnwidth - 4\tabcolsep) * \real{0.30}}
  >{\raggedright\arraybackslash}p{(\columnwidth - 4\tabcolsep) * \real{0.26}}
  >{\raggedright\arraybackslash}p{(\columnwidth - 4\tabcolsep) * \real{0.44}}@{}}
\toprule\noalign{}
\textbf{Riesgo} & \textbf{Detección} & \textbf{Respuesta prevista} \\
\midrule\noalign{}
\endhead
\bottomrule\noalign{}
\endlastfoot
Indisponibilidad o agotamiento de cuota del servicio de cartografía & Error en la respuesta o elementos sin resolver en la matriz & Calcular el recorrido sobre distancias geodésicas, advirtiendo al operador de que la duración es aproximada. \\
Ausencia o corrupción de los modelos entrenados & El servicio no consigue cargar los artefactos al iniciar & Reejecutar el entrenamiento sobre la base vigente y, entretanto, presentar el promedio histórico de cada máquina como estimación de respaldo. \\
Publicación defectuosa de una versión & Errores en el registro de la aplicación o aviso del personal operativo & Redirigir el tráfico a la revisión anterior, que el servicio conserva, y corregir en preproducción. \\
Corrupción de la base de datos & Fallo de las verificaciones de integridad de la carga & Restaurar el respaldo diario o reconstruir la base completa desde el archivo de origen. \\
Pérdida de conectividad del operador en ruta & Sin respuesta desde el dispositivo & Conservar el recorrido calculado en el dispositivo, de modo que la jornada pueda completarse sin nueva consulta. \\
\end{longtable}}

\fuente{Fuente: Elaboración propia}

\subsection{Mantenimiento}\label{mantenimiento}

El sistema no es un producto que se entregue y se dé por terminado, porque su calidad depende de un modelo que envejece a medida que cambia el comportamiento de la red. El mantenimiento previsto se organiza por ello en tres planos: la vigilancia y actualización del modelo, los cambios anticipados sobre el sistema construido y el horizonte de vida útil estimado.

\subsubsection{Reentrenamiento y vigilancia del modelo}\label{reentrenamiento-y-vigilancia-del-modelo}

El histórico crece de manera sostenida —de cuatrocientas treinta y ocho lecturas activas en 2024 se pasó a setecientas veintiuna en 2025—, de modo que cada mes de operación aporta información que el modelo puede aprovechar. Se prevé un reentrenamiento mensual, viable porque el ciclo completo de reconstrucción de la base y ajuste de los quince modelos, con su búsqueda de hiperparámetros incluida, se completa en menos de diez minutos y puede programarse fuera del horario de uso. Cada ejecución deja un informe fechado con las métricas obtenidas, lo que permite seguir su evolución en el tiempo.

La vigilancia se apoya en tres indicadores. El primero es el error absoluto medio de cada insumo en la validación cruzada, comparado con el de la ejecución anterior y con el del método de referencia; un deterioro sostenido en varias ejecuciones consecutivas indica que el patrón de consumo ha cambiado y exige revisar las variables antes que reentrenar sin más. El segundo es la cobertura empírica de las bandas de predicción, que debe mantenerse próxima al valor declarado y, si se desvía, obliga a recalibrar los percentiles como ya se hizo durante la construcción. El tercero es el contraste periódico entre las existencias predichas y las observadas al llegar a la máquina, que es la verificación que el propio operador puede realizar sin instrumentación alguna.

\subsubsection{Cambios previstos y ciclo de vida}\label{cambios-previstos-y-ciclo-de-vida}

Los cambios anticipados se derivan de las limitaciones identificadas durante la construcción y no de una lista de deseos. En el módulo de enrutamiento, incorporar una fase de mejora local y fraccionar la consulta de la matriz de tiempos son requisitos para que el sistema soporte una red sensiblemente mayor que la actual, y el paso a un enrutamiento de varios vehículos sería necesario si la empresa incorporase un segundo operador en ruta simultánea. En el modelo, añadir un calendario académico y de feriados es la línea con mayor recorrido esperado, dado que varias ubicaciones son hospitalarias o universitarias y su afluencia depende del período lectivo más de lo que el mes por sí solo captura; exige, eso sí, una fuente de datos externa de la que hoy no se dispone. Existe además deuda técnica que conviene saldar antes de la operación continuada: el prototipo carece de un conjunto de pruebas automatizadas, el servicio de aplicación no tiene configurado un analizador estático de código como sí lo tiene la interfaz web, y las estructuras de datos de esta replican manualmente las de aquel, por lo que conviene generarlas de forma automática para que un cambio en una no pueda desincronizarse silenciosamente de la otra.

El ciclo de vida previsto se estima en cuatro a cinco años, plazo tras el cual la evolución de la red y de las herramientas empleadas justificaría una revisión de fondo antes que una nueva corrección. Los dos primeros meses corresponden a un piloto en operación asistida, en paralelo con el esquema actual y contrastando contra la línea base; del tercero al sexto mes se sustituye progresivamente el esquema de rutas fijas, se ajusta la interfaz según el uso real y se salda la deuda técnica; entre el primer y el tercer año el sistema entra en operación estable, con reentrenamiento mensual, vigilancia de indicadores y correcciones menores; del tercero al cuarto se incorporan las mejoras previstas y se adecua el sistema al crecimiento de la red; y hacia el quinto corresponde evaluar de fondo la arquitectura y el enfoque de modelado frente a las alternativas disponibles entonces.

En conjunto, este capítulo ha mostrado que la propuesta es construible con tecnologías de uso libre, por un equipo reducido, en ocho semanas y con una inversión modesta, y que su operación posterior cuesta al mes menos que las visitas innecesarias que pretende evitar. Queda por establecer, sin embargo, el criterio con el que se decidirá si el sistema efectivamente genera el valor esperado para Cafetalino. El siguiente capítulo aborda esa cuestión, definiendo el diseño experimental y los criterios de validación del producto mínimo viable frente a la operación actual de la empresa.

% \section{Validación y diseño experimental}\label{validacion-y-diseno-experimental}

El capítulo anterior cerró con una pregunta abierta: construido el sistema y previsto su despliegue, falta el criterio con el que se decidirá si genera realmente el valor esperado. Responderla exige convertir los criterios de éxito enunciados en la etapa Definir y las hipótesis de valor formuladas bajo el enfoque LEAN en algo verificable, es decir, en indicadores con una definición operativa precisa, un instrumento de registro y un umbral de aprobación. Este capítulo delimita primero qué se somete a validación y con qué diseño experimental, y establece después la regla con la que se declarará el producto mínimo viable validado, ajustable o refutado.

\subsection{Alcance de la validación y diseño del piloto}\label{alcance-de-la-validacion-y-diseno-del-piloto}

Lo que se valida es el producto mínimo viable definido en el capítulo de metodología —el modelo predictivo de consumo, el módulo de optimización de recorridos y la interfaz de apoyo a la decisión—, operando sobre las diecisiete ubicaciones activas de la red. Queda fuera lo que el alcance excluyó desde el inicio: al no instalarse sensores ni telemetría, el sistema no mide el nivel real de existencias sino que lo infiere del histórico de reposiciones. La distinción es determinante, porque lo que se pone a prueba no es la exactitud de una medición física sino la calidad de una decisión logística tomada sin instrumentar las máquinas, que es precisamente la apuesta del proyecto. La validación fuera de línea, ya ejecutada y reportada en la Tabla~\ref{tab:desempeno-de-los-modelos}, responde a si el modelo aprende algo aprovechable; la validación operativa que aquí se diseña responde a una pregunta más exigente, la de si esa capacidad se traduce en menos visitas inútiles y menos desabastecimientos, y es la que sustenta el criterio final.

El contraste se realiza mediante un piloto de ocho semanas en operación asistida, coherente con el periodo de medición comprometido en la Tabla~\ref{tab:criterios-de-exito} y con la primera fase del ciclo de vida descrito en el capítulo anterior. Durante ese lapso el sistema opera en paralelo con el procedimiento vigente y no lo sustituye: cada jornada emite su recomendación —las máquinas que a su juicio requieren recarga y el recorrido propuesto para atenderlas— mientras el operador ejecuta su esquema habitual de rutas por zonas, y ambas decisiones quedan registradas junto con lo efectivamente repuesto. La comparación en paralelo se prefiere al reemplazo directo porque no arriesga la continuidad del servicio mientras el sistema no ha demostrado su fiabilidad, y porque evita el sesgo de contrastar dos periodos distintos, en los que la afluencia y la composición de la red no son equiparables.

La unidad de análisis es la visita y cada máquina actúa como su propio control, ya que la recomendación y la decisión del operador se emiten sobre la misma máquina el mismo día; la línea base es el procedimiento actual, con las cifras de referencia recogidas en la etapa Definir. Conviene declarar sin rodeos una limitación del diseño: con diecisiete ubicaciones, el contraste se reporta de forma descriptiva, acompañando cada diferencia observada de su intervalo de confianza, sin perseguir una significación estadística formal que con ese tamaño de red sería difícilmente concluyente. El propósito del piloto es sustentar una decisión de gestión con evidencia cuantificada, no establecer un resultado inferencial.

\subsubsection{Definición operativa de los indicadores}\label{definicion-operativa-de-los-indicadores}

Los criterios de éxito de la Tabla~\ref{tab:criterios-de-exito} enuncian metas pero no dicen cómo se cuenta cada suceso, y sin esa precisión la validación quedaría a merced de la interpretación posterior. La Tabla~\ref{tab:definicion-operativa-de-los-indicadores} cierra esa brecha fijando, para cada indicador, la regla exacta de cómputo y el instrumento con el que se registra durante el piloto, sin repetir las líneas base ni las metas ya establecidas.

\captionof{table}{Definición operativa de los indicadores de validación}\label{tab:definicion-operativa-de-los-indicadores}

{\small\setlength{\tabcolsep}{4pt}
\begin{longtable}[]{@{}
  >{\raggedright\arraybackslash}p{(\columnwidth - 4\tabcolsep) * \real{0.22}}
  >{\raggedright\arraybackslash}p{(\columnwidth - 4\tabcolsep) * \real{0.52}}
  >{\raggedright\arraybackslash}p{(\columnwidth - 4\tabcolsep) * \real{0.26}}@{}}
\toprule\noalign{}
\textbf{Indicador} & \textbf{Definición operativa} & \textbf{Instrumento de registro} \\
\midrule\noalign{}
\endhead
\bottomrule\noalign{}
\endlastfoot
Visita innecesaria & Visita en la que ningún insumo requirió una reposición superior al diez por ciento de la capacidad de su depósito. & Parte de recarga, que ya consigna la cantidad repuesta de cada insumo. \\
Acierto de clasificación & Una máquina requiere recarga si las existencias restantes previstas de algún insumo se agotan antes de la siguiente visita programada; la clasificación emitida se contrasta con lo observado al abrir la máquina. & Comparación entre la recomendación registrada y el parte de recarga de esa visita. \\
Modificación extraordinaria de ruta & Desvío no planificado del recorrido del día originado por un aviso de desabastecimiento. & Parte diario del operador, con el motivo del desvío. \\
Incidencia de desabastecimiento & Reclamo o reembolso atribuido a la ausencia de un insumo en la máquina. & Registro de reclamos de la empresa, contrastado con el año anterior. \\
\end{longtable}}

\fuente{Fuente: Elaboración propia}

La regla de clasificación merece justificación por ser la decisión de mayor incidencia en el resultado. Definir la necesidad de recarga como el agotamiento previsto antes de la próxima visita, y no como un porcentaje fijo de existencias, se corresponde con el horizonte del modelo —que predice el consumo hasta la siguiente visita y no en un plazo fijo de días— y evita penalizar por igual a máquinas de rotación muy distinta. Junto a la exactitud global se vigila por separado la tasa de falsos negativos, esto es, la proporción de máquinas declaradas sin necesidad de recarga que sin embargo se agotan, porque los dos errores no cuestan lo mismo: un falso positivo produce una visita evitable, mientras que un falso negativo produce una máquina vacía, que es justamente el problema que el trabajo pretende resolver.

\subsubsection{Amenazas a la validez}\label{amenazas-a-la-validez}

Cuatro factores condicionan la interpretación de los resultados y conviene enunciarlos antes de obtenerlos \citep{shadish2002}. El histórico es desigual entre ubicaciones, de modo que las máquinas incorporadas recientemente tendrán previsiblemente un desempeño inferior y el resultado agregado debe leerse también desagregado por antigüedad. El operador sabe que su trabajo está siendo contrastado, lo que puede alterar su conducta y aproximarla artificialmente a la del sistema. Ocho semanas no cubren el ciclo anual completo, por lo que los indicadores de periodo anual solo admitirán en ese plazo una lectura de tendencia. Y la cobertura empírica de las bandas de predicción se sitúa en torno al ochenta y tres por ciento frente al noventa nominal, según se documentó en el capítulo anterior, de manera que los intervalos deben presentarse al operador como orientación y no como garantía.

\subsection{Criterio final de validación}\label{criterio-final-de-validacion}

Definidos los indicadores y el protocolo, resta establecer la regla con la que se emitirá el veredicto al término del piloto. Se formula como un criterio compuesto y decidido de antemano, de modo que el resultado no dependa de la interpretación posterior de las cifras, y se acompaña de los tres desenlaces que el ciclo Construir--Medir--Aprender contempla \citep{ries2011}.

El producto mínimo viable se declarará validado si, al término de las ocho semanas, se cumplen simultáneamente dos condiciones necesarias y una de no regresión. La primera condición necesaria es que la clasificación de máquinas alcance una exactitud igual o superior al 85 \%, meta ya comprometida en los criterios de éxito, con una tasa de falsos negativos que no supere el 10 \% de las máquinas clasificadas como sin necesidad de recarga. La segunda es que la proporción de visitas innecesarias sobre el conjunto de recomendaciones se sitúe en el 5 \% o por debajo, frente al 10 \% de la línea base. La condición de no regresión exige que ni las modificaciones extraordinarias de ruta ni las incidencias de desabastecimiento se deterioren respecto del mismo periodo del año anterior.

La asimetría entre unas condiciones y otras responde al periodo de medición de cada indicador. La exactitud de clasificación y la tasa de visitas innecesarias se observan en cada visita, de modo que ocho semanas aportan observaciones suficientes para pronunciarse. En cambio, las metas de reducción del 50 \% en modificaciones de ruta y del 30 \% en incidencias están definidas sobre doce meses y frente a bases anuales pequeñas —alrededor de diez eventos al año en el primer caso—, por lo que exigir su cumplimiento en ocho semanas sería estadísticamente ingenuo; durante el piloto operan como salvaguardas frente a un empeoramiento y su verificación plena se pospone a la operación estable.

Del veredicto se siguen tres desenlaces. Perseverar corresponde al cumplimiento de la regla anterior: el sistema pasa a la fase de adopción prevista en el ciclo de vida, sustituyendo progresivamente el esquema de rutas fijas bajo la vigilancia de indicadores descrita en el plan de mantenimiento. Ajustar se activa cuando la exactitud queda entre el 75 \% y el 85 \%, o cuando alcanza la meta pero la tasa de falsos negativos la excede: la propuesta se mantiene y la intervención se dirige a reentrenar el modelo con los datos del propio piloto, revisar el umbral de agotamiento que separa ambas clases y recalibrar las bandas si procede, tras lo cual se ejecuta un segundo ciclo de medición de cuatro semanas.

Pivotar se reserva para el caso en que la exactitud quede por debajo del 75 \%, o en que el sistema no mejore de forma apreciable al método de referencia constituido por el promedio histórico de consumo de cada máquina. Ese resultado no sería un fallo de implementación sino la refutación de la hipótesis central del trabajo: que el registro histórico de reposiciones contiene información suficiente para anticipar la demanda sin instrumentar las máquinas. La propuesta debería entonces reorientarse hacia la captura directa de datos de consumo, aceptando el costo de hardware que el alcance actual excluye deliberadamente. Enunciar ese umbral por anticipado es lo que hace del diseño una validación y no una justificación, porque somete la hipótesis a una prueba que puede reprobar.

Con ello queda establecido el criterio con el que se juzgará la propuesta, cerrando el recorrido que va del problema identificado a la solución construida y de esta a la evidencia con la que se decidirá su adopción. El capítulo siguiente recoge las conclusiones del trabajo, contrastando lo alcanzado con los objetivos planteados al inicio, y señala las líneas de desarrollo que la propuesta deja abiertas.

% \section{Conclusiones y trabajo futuro}\label{conclusiones-y-trabajo-futuro}

El trabajo partió de un problema verificable en la operación de Cafetalino: un esquema de reabastecimiento apoyado en calendarios fijos y en la experiencia del operador, que produce a la vez visitas a máquinas que aún tienen insumos y desabastecimientos en las que se agotan antes de lo previsto. La propuesta lo abordó anticipando el consumo de cada emplazamiento a partir del registro histórico de reposiciones que la empresa ya lleva, convirtiendo esa proyección en una lista de máquinas que requieren atención y ordenando su visita sobre tiempos de conducción reales. Su validez descansa en que no exige instrumentar las máquinas, cuesta al mes menos que las visitas evitables que suprime y deja la decisión final en el operador, a quien informa en lugar de sustituir. Este capítulo contrasta lo alcanzado con los objetivos planteados, delimita la aportación y sus límites, y señala las líneas que el trabajo deja abiertas.

\subsection{Cumplimiento de los objetivos}\label{cumplimiento-de-los-objetivos}

El objetivo general perseguía mejorar la gestión del reabastecimiento reduciendo las visitas innecesarias y los desabastecimientos mediante un sistema de predicción de la demanda y planificación dinámica de rutas. Los cinco objetivos específicos que lo desagregan se contrastan en la Tabla~\ref{tab:cumplimiento-de-los-objetivos}, con el resultado obtenido y el lugar del documento donde reside la evidencia que lo respalda, de modo que el cumplimiento pueda verificarse y no solo afirmarse.

\captionof{table}{Contraste de los objetivos específicos con los resultados obtenidos}\label{tab:cumplimiento-de-los-objetivos}

{\small\setlength{\tabcolsep}{4pt}
\begin{longtable}[]{@{}
  >{\raggedright\arraybackslash}p{(\columnwidth - 4\tabcolsep) * \real{0.26}}
  >{\raggedright\arraybackslash}p{(\columnwidth - 4\tabcolsep) * \real{0.50}}
  >{\raggedright\arraybackslash}p{(\columnwidth - 4\tabcolsep) * \real{0.24}}@{}}
\toprule\noalign{}
\textbf{Objetivo específico} & \textbf{Resultado alcanzado} & \textbf{Evidencia} \\
\midrule\noalign{}
\endhead
\bottomrule\noalign{}
\endlastfoot
Caracterizar los patrones históricos de consumo y reabastecimiento. & Alcanzado. Los treinta y seis identificadores de equipo del archivo se consolidaron en diecisiete ubicaciones canónicas por coordenada, y de las lecturas activas quedaron 2.439 observaciones útiles tras descartar las que carecían de historial previo. & Apartado~\ref{origen-y-preparacion-de-los-datos} \\
Desarrollar un módulo de estimación de reposición sin hardware adicional. & Alcanzado. Cinco modelos independientes, uno por insumo, acompañados de bandas de incertidumbre cuya cobertura empírica se verificó sobre pliegues no vistos en lugar de declararse. & Tabla~\ref{tab:desempeno-de-los-modelos} \\
Desarrollar un módulo de optimización de recorridos dinámicos. & Alcanzado. El recorrido se construye sobre la matriz de tiempos de conducción reales, parte de la posición del operador y suma el tiempo de servicio de cada máquina. & Apartado~\ref{modulo-de-enrutamiento-dinamico} \\
Desarrollar una interfaz web de apoyo a la decisión. & Alcanzado. La aplicación organiza la jornada en tres columnas y muestra en cada tarjeta el intervalo además de la estimación puntual, para que el operador reserve su criterio ante bandas anchas. & Apartado~\ref{interfaz-web-de-apoyo-a-la-decision} \\
Evaluar el desempeño frente al esquema de rutas fijas actual. & Parcialmente alcanzado. La evaluación fuera de línea se ejecutó y sus métricas están reportadas; el contraste operativo contra el esquema vigente quedó diseñado con umbrales decididos de antemano, pero no ejecutado. & Apartado~\ref{criterio-final-de-validacion} \\
\end{longtable}}

\fuente{Fuente: Elaboración propia}

El quinto objetivo merece discusión explícita por ser el único que no se cierra. Su cumplimiento pleno exige el piloto de ocho semanas comprometido en los criterios de éxito de la Tabla~\ref{tab:criterios-de-exito}, cuyo periodo excede el horizonte del trabajo académico y depende de la disponibilidad del personal de la empresa. Sí quedó resuelta la condición que hace posible esa evaluación: los indicadores tienen ahora definición operativa e instrumento de registro, y la regla de aprobación fija umbrales anteriores a la obtención de los datos, incluido el que refutaría la hipótesis central del trabajo. La diferencia no es menor, porque un criterio establecido de antemano impide que el resultado se acomode después a la interpretación más favorable.

\subsection{Aportaciones y limitaciones}\label{aportaciones-y-limitaciones}

La aportación no reside en los algoritmos empleados, que son estándar, sino en mostrar que un registro administrativo que la empresa ya lleva por obligación contable contiene señal suficiente para decidir la logística de reposición sin instrumentar las máquinas. El estado del arte revisado atiende mayoritariamente a redes que disponen de telemetría, condición que la pequeña y mediana empresa del sector en el contexto local no cumple ni puede costear en el corto plazo, de modo que la brecha cubierta es la de operar con el dato que existe en lugar del que sería deseable. A ello se suma el registro documentado de lo que no funcionó —tres variantes de variables descartadas y la recalibración de unas bandas cuya cobertura nominal no se sostenía empíricamente—, que delimita hasta dónde llega el enfoque y ahorra ese recorrido a quien lo retome.

Las limitaciones deben enunciarse con la misma claridad. El coeficiente de determinación obtenido, entre 0,27 y 0,37, indica que parte sustancial de la variabilidad del consumo permanece sin explicar, por una causa informativa antes que algorítmica: el histórico registra cuándo y cuánto se repuso, pero no la afluencia de personas ni los factores que la determinan. La red es además pequeña, diecisiete ubicaciones, lo que obliga a reportar cualquier contraste de forma descriptiva y limita la generalización de las cifras a redes de otro tamaño. Por último, el recorrido que entrega el optimizador es una construcción de buena calidad y no un óptimo demostrado, decisión deliberada para privilegiar la respuesta inmediata, cuya diferencia práctica es reducida hoy pero dejaría de serlo si la red creciera.

\subsection{Líneas de trabajo futuro}\label{lineas-de-trabajo-futuro}

Los cambios técnicos previstos sobre el sistema construido, la deuda técnica pendiente y el ciclo de vida estimado se detallaron ya en el apartado~\ref{cambios-previstos-y-ciclo-de-vida} y no se reiteran aquí. Las líneas que siguen se sitúan en el plano de las perspectivas que el trabajo abre más allá del caso de Cafetalino, ordenadas por su capacidad de superar las limitaciones recién enunciadas.

La primera es la incorporación de fuentes externas de demanda. Un calendario académico y de feriados, la condición climática y la agenda de eventos de cada zona aportarían la información de afluencia que el histórico no contiene, y son por tanto la vía con mayor recorrido esperado para elevar el techo predictivo identificado. Su efecto no sería solo de exactitud, ya que convertiría un modelo reactivo, que extrapola el comportamiento pasado de cada emplazamiento, en uno capaz de anticipar cambios de régimen como el inicio de un periodo lectivo o un feriado prolongado.

La segunda es una estrategia híbrida de instrumentación selectiva. En lugar de la disyuntiva entre instrumentar toda la red o ninguna máquina, cabe dotar de telemetría solo a los emplazamientos de mayor rotación y emplear su lectura directa para corregir la predicción del resto, que comparte con ellos calendario y patrón de afluencia. Esa vía preserva la ventaja económica que da sentido a la propuesta y permite cuantificar el error real del enfoque sin instrumentación, que hoy solo puede estimarse indirectamente.

La tercera es la transferibilidad. Lo que el sistema aprovecha no es una propiedad de las bebidas calientes sino la estructura del dato —visita, intervalo transcurrido y cantidad repuesta—, común a cualquier operación de reposición periódica que documente sus visitas, de modo que la aportación se proyecta a otras redes de expendio automático y, más ampliamente, a la reposición de puntos de venta menudos, al mantenimiento de equipamiento distribuido o a la lectura de medidores. La cuarta es el escalado logístico: el enrutamiento de varios vehículos con ventanas horarias y restricciones de capacidad, que dejaría de ser una mejora opcional para volverse un requisito si la red crece o si se incorpora un segundo operador en ruta simultánea.

En síntesis, el trabajo deja demostrado que la propuesta es construible con tecnologías de uso libre, por un equipo reducido y con una inversión modesta, y que el histórico disponible contiene señal aprovechable por encima del promedio simple de cada máquina. Queda por demostrar que esa señal basta para mejorar la operación real, cuestión que solo el piloto resuelve y que el criterio del capítulo anterior permite decidir sin ambigüedad. Haber enunciado por anticipado qué resultado obligaría a abandonar el enfoque es lo que convierte la propuesta en una innovación verificable y no en una expectativa.


% ---- Bibliografía ------------------------------------------------
\imprimirbibliografia

% ---- Anexos ------------------------------------------------------
\appendix
\renewcommand{\thesection}{Anexo \Alph{section}}
% los anexos necesitan más hueco para su numeración en el índice
\addtocontents{toc}{\protect\setlength{\protect\cftsecnumwidth}{5.2em}}
\titleformat{\section}{\titlefont\fontsize{18}{22}\selectfont\color{unirazul}}{\thesection.}{0.5em}{}
\section{Entrevista al Gerente de Cafetalino}\label{anexo-entrevista-al-gerente-de-cafetalino}

\textbf{GUÍA PARA LA ENTREVISTA AL GERENTE PROPIETARIO DE CAFETALINO}

El objetivo de la entrevista es conocer y comprender la empresa y sus procedimientos a fin de identificar problemas que podrían encararse con la ayuda de sistemas basados en inteligencia artificial.

\begin{enumerate}
\def\labelenumi{\arabic{enumi}.}
\item
  ¿Tiene identificado algún problema recurrente, proceso pesado, o algo que le gustaría optimizar? Por ejemplo, predecir qué meses se venderá más o automatizar alguna tarea.
\item
  ¿Guarda algún tipo de datos o registros de sus operaciones?
\item
  ¿Considera que una aplicación de inteligencia artificial aportaría valor a sus operaciones?
\item
  ¿Estaría dispuesto a compartirla para implementar soluciones de inteligencia artificial con el objetivo de mejorar sus operaciones?
\item
  ¿Alguna vez intento instalar sensores para hacer seguimiento al consumo de sus insumos?
\item
  ¿Alguna vez ha tenido algún incidente con sus clientes? ¿De qué tipo?
\item
  ¿Cómo ha resuelto estos problemas?
\item
  ¿Cuál es el procedimiento de recarga de sus máquinas dispensadoras?
\item
  ¿Alguna vez ha tenido algún problema con este proceso? ¿De qué tipo?
\end{enumerate}

\textbf{Desarrollo de la Entrevista al Gerente Propietario de Cafetalino}

El objetivo de la entrevista es conocer y comprender la empresa y sus procedimientos a fin de identificar problemas que podrían encararse con la ayuda de sistemas basados en inteligencia artificial.

\begin{enumerate}
\def\labelenumi{\arabic{enumi}.}
\item
  ¿Tiene identificado algún problema recurrente, proceso pesado, o algo que le gustaría optimizar? Por ejemplo, predecir qué meses se venderá más o automatizar alguna tarea.
\end{enumerate}

\begin{quote}
\emph{Tenemos muchas cosas que queremos incorporar y mejorar a lo que actualmente hacemos:}
\end{quote}

\begin{enumerate}
\def\labelenumi{\arabic{enumi}.}
\item
  \emph{Enviar notificaciones a los clientes de que su tarjeta esta con crédito insuficiente y que les sugiera recargar.}
\item
  \emph{Nuestro sistema lleva registro de todas las operaciones con tarjeta. Nos gustaría que también registre las operaciones con monedas y que esto este diferenciado: Operaciones con tarjeta y operaciones con monedas.}
\item
  \emph{Quisiéramos abrir nuestro sistema a los clientes mediante una app y/o una web para consultas de saldo, recargas de tarjeta u otros beneficios}
\item
  \emph{Nos gustaría poder conocer, de todas las recargas a las tarjetas, cuantas de ellas ya se han consumido y cuánto es el saldo, que se constituye en producto por entregar, que ha sido pagado por adelantado. Así hay muchas cosas, necesitamos muchas cosas más, hay muchas ideas.}
\item
  \emph{Por ejemplo, un sistema de inventario de insumos, que son distribuidos a las maquinas, estos se convierten en productos y desde las maquinas son comercializados. Que me devuelva saldos para verificarlos físicamente o que pueda tener una tendencia para pronosticar los requerimientos de recarga de insumos a las maquinas.}
\end{enumerate}

\begin{quote}
Esta quinta idea me parece interesante y viable para darle una solución con inteligencia artificial, que es lo que me gustaría aplicar.
\end{quote}

\begin{enumerate}
\def\labelenumi{\arabic{enumi}.}
\setcounter{enumi}{1}
\item
  ¿Guarda algún tipo de datos o registros de sus operaciones?
\end{enumerate}

\begin{quote}
\emph{Si, tenemos datos de toda la actividad de Cafetalino desde el 2018, de hecho, también se podría estimar la cantidad de insumos que hemos usado en función a las ventas.}

\emph{Hay un registro de todas las operaciones con tarjeta. Nos gustaría que también registre las operaciones con monedas y que esto este diferenciado: Operaciones con tarjeta y operaciones con monedas.}
\end{quote}

\begin{enumerate}
\def\labelenumi{\arabic{enumi}.}
\setcounter{enumi}{2}
\item
  ¿Considera que una aplicación de inteligencia artificial aportaría valor a sus operaciones?
\end{enumerate}

\begin{quote}
\emph{Por supuesto, además sería una forma de modernizar nuestras operaciones, de estar a la vanguardia}
\end{quote}

\begin{enumerate}
\def\labelenumi{\arabic{enumi}.}
\setcounter{enumi}{3}
\item
  ¿Estaría dispuesto a compartir su información para implementar soluciones de inteligencia artificial con el objetivo de mejorar sus operaciones?
\end{enumerate}

\begin{quote}
\emph{Mmm, si, para mejorar nuestras operaciones por supuesto que sí, pero la condición es manejar con mucha discreción, no quisiéramos que, de ninguna manera, se filtre la información y cosas que nos ha costado aprender sean aprovechadas por otros para hacernos la competencia. De hecho, tendríamos que hacer un contrato de confidencialidad.}

Por supuesto, no hay de que preocuparse pues se manejaría con mucha profesionalidad.
\end{quote}

\begin{enumerate}
\def\labelenumi{\arabic{enumi}.}
\setcounter{enumi}{4}
\item
  ¿Alguna vez intento instalar sensores para hacer seguimiento al consumo de sus insumos?
\end{enumerate}

\begin{quote}
\emph{Si, en una primera instancia se realizó una prueba piloto y se instalaron sensores en algunas máquinas para saber si era necesario recargar vasos, y que el sistema nos alerte, vía mensaje, cuando quedaba una pequeña cantidad. Sin embargo, se tuvieron problemas de conectividad y no prospero la idea, además se limitaba solo a los vasos y no a los demás insumos.}

\emph{Se exploro la posibilidad de instalar más sensores, uno para cada uno de insumos, pero el presupuesto inicial era muy alto, no había la certeza de que funcionara, pues los problemas de conectividad persistían, además había limitaciones técnicas especialmente con los insumos sólidos, yyy también con los líquidos.}
\end{quote}

\begin{enumerate}
\def\labelenumi{\arabic{enumi}.}
\setcounter{enumi}{5}
\item
  ¿Alguna vez ha tenido algún incidente con sus clientes? ¿De qué tipo?
\end{enumerate}

\begin{quote}
\emph{Si. Alguna vez cuando una maquina se ha quedado sin insumos o quedó muy poca cantidad de alguno de ellos, ocurrió que: o no entrego el producto o este no tenía la calidad a la que el cliente estaba acostumbrado. Esto genera, de forma inmediata sentimientos de frustración en el cliente, decepción e impotencia, no saben a quién reclamar, y una sensación de pérdida o robo, pierden la confianza y muchos ya no se animan a volver a interactuar con la máquina.}

\emph{En algunos casos estos sentimientos escalaron en el nivel de violencia, algunos usuarios con menor tolerancia, tuvieron comportamientos agresivos, golpearon los aparatos, y en ocasiones provocaron daños a nuestras maquinas; en otras ocasiones llamaron al número de contacto y nos maltrataron, nos acusaron de robo y no aceptaron nuestras disculpas, ni nos dejaron darles ninguna explicación, hasta nos amenazaron con procesos judiciales, nos difamaron en las redes sociales y en su entorno de trabajo. Alguna vez hasta tuvimos ir a su entorno de trabajo, les explicamos lo sucedido y a modo de disculpa les ofrecimos bebidas a todos los trabajadores.}

\emph{En un caso en particular, recibimos la denuncia de sus compañeros de trabajo, que el sujeto provocaba el fallo de forma intencionada, introduciendo otros objetos o monedas de otros países y que luego protestaba en voz alta. Bueno, también nos comentaron que esa era su forma cotidiana de actuar y hasta tuvimos que retirar nuestro equipo de esa ubicación para evitar el daño, a los equipos y a la imagen de la empresa.}
\end{quote}

\begin{enumerate}
\def\labelenumi{\arabic{enumi}.}
\setcounter{enumi}{6}
\item
  ¿Cómo ha resuelto estos problemas?
\end{enumerate}

\begin{quote}
\emph{Bueno, en estas situaciones, los que nos conocen nos llamaban y nos comunicaban lo que sucedió; nosotros le devolvíamos su dinero y le dábamos de cortesía la bebida, cuando era posible.}

\emph{Luego pusimos un número de teléfono de contacto en todas nuestras ubicaciones, para que cualquier cliente, aunque no nos conozca, pueda reportar el suceso y procedíamos de la misma manera, devolverle su dinero y si era posible darle su bebida de cortesía}

\emph{Sin embargo, como obviamente transcurría un tiempo para que podamos ir a recargar la maquina y ponerla nuevamente en servicio y no siempre era posible darle la bebida de cortesía. Con la devolución del dinero, no hubo problema, inicialmente hablamos con vecinos de nuestras ubicaciones que nos hacían el favor de devolver el dinero y luego nosotros les reponíamos, pero ahora con la facilidad de las trasferencias por Qr eso se ha facilitado enormemente}.
\end{quote}

\begin{enumerate}
\def\labelenumi{\arabic{enumi}.}
\setcounter{enumi}{7}
\item
  ¿Cuál es el procedimiento de recarga de sus máquinas dispensadoras?
\end{enumerate}

\begin{quote}
\emph{Al principio cuando teníamos pocas maquinas, y nos estábamos haciendo conocer, hacíamos el recorrido por una ruta fija, revisando todas nuestras máquinas y recargando en función a lo consumido. Esto lo hacíamos una vez por semana, una vez cada 2 semanas, una vez cada 10 días, fuimos ajustando la frecuencia en base a la cantidad consumida.}

\emph{Después conforme nos fuimos haciendo conocer, aumentamos la frecuencia de recarga y la cantidad de máquinas, luego agrupamos las maquinas en función a su ubicación y aumentamos el número de rutas. Actualmente, revisamos dos rutas cada día, y tenemos personal que se ocupa de ello.}
\end{quote}

\begin{enumerate}
\def\labelenumi{\arabic{enumi}.}
\setcounter{enumi}{8}
\item
  ¿Alguna vez ha tenido algún problema con este proceso? ¿De qué tipo?
\end{enumerate}

\begin{quote}
\emph{Bueno, si se le puede llamar problema, a veces ocurre que, al revisar una máquina, no ha tenido el consumo esperado y no requiere recarga de insumos.}

\emph{En otras ocasiones, nos hemos visto en la necesidad de modificar nuestra ruta porque se nos reportó que una máquina, de la misma ruta se quedó sin insumos y dejo de prestar el servicio, de manera que es necesario priorizarla, o cuando se trata de una máquina de otra ruta, es necesario priorizarla e incluirla en la secuencia. Esto nos obliga a llevar una cantidad extra de insumos. En otras ocasiones el reporte de una maquina fuera de servicio, llegó al terminar el recorrido y se tuvo que volver a realizar la rutina por una sola máquina: preparar paquetes de recarga de insumos, dirigirnos a la ubicación lo más rápido posible, revisar la máquina, hacerle un mantenimiento rutinario, reabastecerla y ponerla en servicio.}

¿Tiene idea de cuántas veces ocurren estos eventos a la semana, al mes o al año?

\emph{Maquinas que no tienen necesidad de recarga aproximadamente 1 de cada 10.}

\emph{Eventos de modificación de la ruta de reabastecimiento, 10 al año.}

Se me ocurre una aplicación que permita que procesar el historial de recargas (2018-2026) para predecir el desabastecimiento y planificar automáticamente rutas dinámicas óptimas.

\emph{Perfecto, me parece interesante la idea.}
\end{quote}

\section{Entrevista a un Operador de Cafetalino}\label{anexo-entrevista-a-un-operador-de-cafetalino}

\textbf{GUÍA PARA LA ENTREVISTA A UN OPERADOR DE RECARGA DE LA RED DE MÁQUINAS DISPENSADORAS DE CAFETALINO}

El objetivo de la entrevista es conocer y comprender el proceso de recarga de insumos a la red de máquinas dispensadoras de Cafetalino, los problemas y particularidades del mismo.

\begin{enumerate}
\def\labelenumi{\arabic{enumi}.}
\item
  ¿Cuál es el procedimiento de recarga de las máquinas dispensadoras?
\item
  ¿Cómo se definen las rutas de recarga?
\item
  ¿Alguna vez ha tenido que modificar la ruta de recarga? ¿Por qué?
\item
  ¿Cuán eficiente considera que es este procedimiento?
\item
  ¿Consideraría útil un sistema que le ayude a priorizar sus rutas en función a la demanda de recarga de las máquinas de la red de Cafetalino?
\end{enumerate}

\textbf{DESARROLLO DE LA ENTREVISTA A UN OPERADOR DE RECARGA DE LA RED DE MÁQUINAS DISPENSADORAS DE CAFETALINO}

\begin{enumerate}
\def\labelenumi{\arabic{enumi}.}
\item
  ¿Cuál es el procedimiento de recarga de las máquinas dispensadoras?
\end{enumerate}

\begin{quote}
\emph{Cada día recorremos dos rutas que comprenden entre 4 y 6 maquinas, una en la mañana y otra en la tarde.}

\emph{Antes de salir preparamos los paquetes de recarga de insumos, uno por maquina y 2 extra por si es necesario. Luego hacemos el recorrido y en cada ubicación, apagamos la máquina, la revisamos, le hacemos limpieza y un mantenimiento rutinario, la recargamos y la encendemos nuevamente para ponerla en servicio, verificando su correcto funcionamiento, que no haya fugas, ruidos extraños, ni alertas de ningún tipo.}
\end{quote}

\begin{enumerate}
\def\labelenumi{\arabic{enumi}.}
\setcounter{enumi}{1}
\item
  ¿Cómo se definen las rutas de recarga?
\end{enumerate}

\begin{quote}
\emph{Las rutas están predefinidas, cada ruta cubre un grupo de máquinas que están más o menos por la misma zona, el número de máquinas está en función a la distancia entre maquinas.}
\end{quote}

\begin{enumerate}
\def\labelenumi{\arabic{enumi}.}
\setcounter{enumi}{2}
\item
  ¿Alguna vez ha tenido que modificar la ruta de recarga? ¿Por qué?
\end{enumerate}

\begin{quote}
\emph{Si, a veces, nos avisan que hay una máquina fuera de servicio, si está en la misma ruta, le damos prioridad y vamos primero a esa ubicación, y luego cambiamos la secuencia para no estar yendo de un lado a otro, lo grave es si es de otra ruta, porque siempre hay que priorizar la que se quedó fuera de servicio y luego volver a la ruta normal o invertir la secuencia para no estar yendo de un lado a otro, para estos casos es que preparamos 2 paquetes de recarga extra.}
\end{quote}

\begin{enumerate}
\def\labelenumi{\arabic{enumi}.}
\setcounter{enumi}{3}
\item
  ¿Cuán eficiente considera que es este procedimiento?
\end{enumerate}

\begin{quote}
\emph{No es muy eficiente, porque como le dije a veces es necesario cambiar la secuencia de la ruta o hasta ir a otras rutas y muchas veces hay que hacer horas extra para terminar de revisar todas las maquinas. Además, a veces llegamos a una ubicación que por algún motivo no ha tenido el consumo regular, porque hubo corte de energía eléctrica en la zona o alguna otra cosa, lo que pasa en estos casos es que no requiere recarga y hacerle el mantenimiento se complica porque los depósitos están llenos, y bueno solo se la limpia y se la pone en servicio nuevamente, pero es solo una pérdida de tiempo}
\end{quote}

\begin{enumerate}
\def\labelenumi{\arabic{enumi}.}
\setcounter{enumi}{4}
\item
  ¿Consideraría útil un sistema que le ayude a priorizar sus rutas en función a la predicción de la demanda de recarga de las máquinas de la red de Cafetalino?
\end{enumerate}

\begin{quote}
\emph{Uhh, sería una gran cosa, pues se evitarían estas pérdidas de tiempo, la ruta se programaría con las maquinas que realmente requieren recarga, y no se tendría que modificar las rutas a medio recorrido, o quizás sí, pero sería excepcionalmente.}
\end{quote}

\section{Encuesta a los Clientes de Cafetalino}\label{anexo-encuesta-a-los-clientes-de-cafetalino}

\textbf{CUESTIONARIO DE LA ENCUESTA A LOS CLIENTES DE CAFETALINO}

Esta encuesta es anónima y tiene el objetivo de recabar información en relación a la perspectiva y sentimiento de los clientes de Cafetalino frente a fallos en el servicio.

Por favor responda con la mayor sinceridad posible.

\begin{enumerate}
\def\labelenumi{\arabic{enumi}.}
\item
  Indique con qué frecuencia usa los servicios de Cafetalino.
\end{enumerate}

\begin{itemize}[label=$\square$,itemsep=1pt,leftmargin=3em]
\item Mas de 2 veces/semana
\item 1 vez/semana
\item 1 vez/mes
\item 2 veces/semana
\item 2 veces/mes
\item Menos de 1 vez/mes
\end{itemize}

\begin{enumerate}
\def\labelenumi{\arabic{enumi}.}
\setcounter{enumi}{1}
\item
  Exprese su grado de satisfacción con el servicio que le brinda Cafetalino, en una escala de 1 a 5: donde: \textbf{5} = Muy bueno, \textbf{4} = Bueno, \textbf{3} = Regular, \textbf{2} = Malo y \textbf{1} = Muy malo.
\end{enumerate}

\begin{itemize}[label=$\square$,itemsep=1pt,leftmargin=3em]
\item Grado de satisfacción
\end{itemize}

\begin{enumerate}
\def\labelenumi{\arabic{enumi}.}
\setcounter{enumi}{2}
\item
  Al usar la máquina de Cafetalino, alguna vez le ocurrió que (puede marcar varias opciones):
\end{enumerate}

\begin{itemize}[label=$\square$,itemsep=1pt,leftmargin=3em]
\item No le entrego el Producto
\item El producto no tenía la calidad que normalmente tiene
\item Nunca tuve ese tipo de experiencia
\end{itemize}

Si su respuesta a la pregunta fue la opción c), le agradecemos su ayuda.

Si su respuesta a la pregunta fue la opción a) o b), o ambas por favor responda las siguientes preguntas.

\begin{enumerate}
\def\labelenumi{\arabic{enumi}.}
\setcounter{enumi}{3}
\item
  Indique cuantas veces el servicio de Cafetalino no fue de su agrado.
\end{enumerate}

\begin{itemize}[label=$\square$,itemsep=1pt,leftmargin=3em]
\item Una vez
\item Dos veces
\item Mas de 2 veces
\end{itemize}

\begin{enumerate}
\def\labelenumi{\arabic{enumi}.}
\setcounter{enumi}{4}
\item
  Marque que sentimientos le genero este servicio inadecuado (Puede marcar varias opciones).
\end{enumerate}

\begin{itemize}[label=$\square$,itemsep=1pt,leftmargin=3em]
\item Decepción
\item Impotencia
\item Desconfianza
\item Frustración
\item Pérdida o Robo
\item Confianza
\end{itemize}

\begin{enumerate}
\def\labelenumi{\arabic{enumi}.}
\setcounter{enumi}{5}
\item
  Indique si reporto el suceso a Cafetalino.
\end{enumerate}

\begin{itemize}[label=$\square$,itemsep=1pt,leftmargin=3em]
\item SI
\item NO
\end{itemize}

\begin{enumerate}
\def\labelenumi{\arabic{enumi}.}
\setcounter{enumi}{6}
\item
  Indique por qué medio reportó el suceso a Cafetalino.
\end{enumerate}

\begin{itemize}[label=$\square$,itemsep=1pt,leftmargin=3em]
\item Llamé al teléfono de contacto
\item En forma personal
\end{itemize}

\begin{enumerate}
\def\labelenumi{\arabic{enumi}.}
\setcounter{enumi}{7}
\item
  Si reporto el suceso, indique si Cafetalino le dio alguna disculpa y solución (Puede marcar varias opciones).
\end{enumerate}

\begin{itemize}[label=$\square$,itemsep=1pt,leftmargin=3em]
\item Se disculparon
\item Me devolvieron el dinero
\item Me dieron una compensación
\end{itemize}

\begin{enumerate}
\def\labelenumi{\arabic{enumi}.}
\setcounter{enumi}{8}
\item
  Indique cuanto tiempo después volvió a usar el servicio de Cafetalino.
\end{enumerate}

\begin{itemize}[label=$\square$,itemsep=1pt,leftmargin=3em]
\item El mismo día
\item 2 días después
\item 2 semanas después
\item 1 día después
\item 1 semana después
\item Mas de 2 semanas después
\end{itemize}

\begin{quote}
\textbf{OPERACIONALIZACION DE LAS RESPUESTAS A LA ENCUESTA A LOS CLIENTES DE CAFETALINO}
\end{quote}

La siguiente tabla muestra la operacionalización de las respuestas de los clientes de Cafetalino para facilitar su procesamiento y análisis estadístico.

\captionof{table}{Operacionalización de las respuestas de los clientes de Cafetalino}\label{tab:operacionalizacion-de-las-respuestas-de}

\begin{longtable}[]{@{}
  >{\raggedright\arraybackslash}p{(\columnwidth - 4\tabcolsep) * \real{0.0629}}
  >{\raggedright\arraybackslash}p{(\columnwidth - 4\tabcolsep) * \real{0.8428}}
  >{\raggedright\arraybackslash}p{(\columnwidth - 4\tabcolsep) * \real{0.0943}}@{}}
\toprule\noalign{}
\begin{minipage}[b]{\linewidth}\raggedright
\textbf{N.º}
\end{minipage} & \begin{minipage}[b]{\linewidth}\raggedright
\textbf{Pregunta/Respuesta}
\end{minipage} & \begin{minipage}[b]{\linewidth}\raggedright
\textbf{Valor}
\end{minipage} \\
\midrule\noalign{}
\endhead
\bottomrule\noalign{}
\endlastfoot
\textbf{1} & Indique con qué frecuencia usa los servicios de Cafetalino. & \\
& Mas de 2 veces/semana & \textbf{6} \\
& 2 veces/semana & \textbf{5} \\
& 1 vez/semana & \textbf{4} \\
& 2 veces/mes & \textbf{3} \\
& 1 vez/mes & \textbf{2} \\
& Menos de 1 vez/mes & \textbf{1} \\
& No responde & \textbf{0} \\
\textbf{2} & Exprese su grado de satisfacción con el servicio que le brinda Cafetalino, en una escala de 1 a 5. & \\
& Muy bueno & \textbf{5} \\
& Bueno & \textbf{4} \\
& Regular & \textbf{3} \\
& Malo & \textbf{2} \\
& Muy malo & \textbf{1} \\
& No responde & \textbf{0} \\
\textbf{3} & Al usar la máquina de Cafetalino, alguna vez le ocurrió que (puede marcar varias opciones). & \\
& No le entrego el Producto & \textbf{1} \\
& El producto no tenía la calidad que normalmente tiene & \textbf{2} \\
& Nunca tuve ese tipo de experiencia & \textbf{3} \\
& No responde & \textbf{0} \\
\textbf{4} & Indique cuantas veces el servicio de Cafetalino no fue de su agrado. & \\
& Una vez & \textbf{1} \\
& Dos veces & \textbf{2} \\
& Mas de 2 veces & \textbf{3} \\
& No responde & \textbf{0} \\
\textbf{5} & Marque que sentimientos le genero este servicio inadecuado (Puede marcar varias opciones). & \\
& Decepción & \textbf{1} \\
& Frustración & \textbf{2} \\
& Impotencia & \textbf{3} \\
& Pérdida o robo & \textbf{4} \\
& Desconfianza & \textbf{5} \\
& Confianza & \textbf{6} \\
& No responde & \textbf{0} \\
\textbf{6} & Indique si reporto el suceso a Cafetalino. & \\
& SI & \textbf{1} \\
& NO & \textbf{2} \\
& No responde & \textbf{0} \\
\textbf{7} & Indique por qué medio reportó el suceso a Cafetalino. & \\
& Llamé al teléfono de contacto & \textbf{1} \\
& En forma personal & \textbf{2} \\
& No responde & \textbf{0} \\
\textbf{8} & Si reporto el suceso, indique si Cafetalino le dio alguna disculpa y solución (Puede marcar varias opciones). & \\
& Se disculparon & \textbf{1} \\
& Me devolvieron el dinero & \textbf{2} \\
& Me dieron una compensación & \textbf{3} \\
& No responde & \textbf{0} \\
\textbf{9} & Indique cuanto tiempo después volvió a usar el servicio de Cafetalino. & \\
& El mismo día & \textbf{6} \\
& 1 día después & \textbf{5} \\
& 2 días después & \textbf{4} \\
& 1 semana después & \textbf{3} \\
& 2 semanas después & \textbf{2} \\
& Mas de 2 semanas después & \textbf{1} \\
& No responde & \textbf{0} \\
\end{longtable}

\fuente{Fuente: Elaboración Propia}

\textbf{RESULTADOS DE LA ENCUESTA A LOS CLIENTES DE CAFETALINO}

La siguiente tabla muestra los resultados de la encuesta aplicada a los clientes de cafetalino, en cantidad de clientes que dieron una determinada respuesta, de acuerdo con la operacionalización de respuestas detallada en el punto anterior.

\captionof{table}{Resultados de la Encuesta aplicada a los clientes de Cafetalino}\label{tab:resultados-de-la-encuesta-aplicada}

\begin{center}
\begin{tabular}{@{}l*{7}{c}@{}}
\toprule
\multirow{2}{*}{\textbf{Pregunta}} & \multicolumn{7}{c}{\textbf{Respuestas}} \\
\cmidrule(l){2-8}
 & \textbf{0} & \textbf{1} & \textbf{2} & \textbf{3} & \textbf{4} & \textbf{5} & \textbf{6} \\
\midrule
\textbf{1} & 1 & 2 & 3 & 21 & 73 & 95 & 59 \\
\textbf{2} & 2 & 1 & 23 & 15 & 117 & 96 & - \\
\textbf{3} & 2 & 16 & 7 & 229 & - & - & - \\
\textbf{4} & 0 & 19 & 4 & 0 & - & - & - \\
\textbf{5} & 0 & 21 & 18 & 6 & 1 & 10 & 13 \\
\textbf{6} & 1 & 19 & 3 & - & - & - & - \\
\textbf{7} & & 17 & 2 & - & - & - & - \\
\textbf{8} & 2 & 17 & 17 & 11 & - & - & - \\
\textbf{9} & 0 & 1 & 5 & 2 & 9 & 4 & 1 \\
\bottomrule
\end{tabular}
\end{center}

\fuente{Fuente: Elaboración Propia}

La siguiente tabla muestra los porcentajes de clientes que dieron un determinado tipo de respuesta, luego de procesar los resultados anteriores.

\captionof{table}{Resultados de la Encuesta aplicada a los clientes de Cafetalino}\label{tab:resultados-de-la-encuesta-aplicada-2}

\begin{center}
\begin{tabular}{@{}l*{7}{c}@{}}
\toprule
\multirow{2}{*}{\textbf{Pregunta}} & \multicolumn{7}{c}{\textbf{Respuestas}} \\
\cmidrule(l){2-8}
 & \textbf{0} & \textbf{1} & \textbf{2} & \textbf{3} & \textbf{4} & \textbf{5} & \textbf{6} \\
\midrule
\textbf{1} & 0.4 & 0.8 & 1.2 & 8.3 & 28.7 & 37.4 & 23.2 \\
\textbf{2} & 0.8 & 0.4 & 9.1 & 5.9 & 46.1 & 37.8 & -- \\
\textbf{3} & 0.8 & 6.3 & 2.8 & 90.2 & -- & -- & -- \\
\textbf{4} & 0.0 & 82.6 & 17.4 & 0.0 & -- & -- & -- \\
\textbf{5} & 0.0 & 91.3 & 78.3 & 26.1 & 4.4 & 43.5 & 56.5 \\
\textbf{6} & 4.4 & 82.6 & 13.0 & -- & -- & -- & -- \\
\textbf{7} & 0.0 & 89.5 & 10.5 & -- & -- & -- & -- \\
\textbf{8} & 4.3 & 100.0 & 100.0 & 64.7 & -- & -- & -- \\
\textbf{9} & 0.0 & 4.6 & 22.7 & 9.1 & 40.9 & 18.2 & 4.6 \\
\bottomrule
\end{tabular}
\end{center}

\fuente{Fuente: Elaboración Propia}

Las figuras siguientes representan gráficamente los resultados recogidos en las tablas anteriores. Se ordenan según la secuencia del cuestionario, de modo que la primera corresponde a la frecuencia de uso del servicio y la última al comportamiento posterior de los clientes que sufrieron un incidente. Los porcentajes que se citan en el cuerpo del documento proceden de estas mismas respuestas.

\begin{figure}[!htbp]
\centering
\caption{Hábitos de consumo de los clientes de Cafetalino}\label{fig:habitos-de-consumo-de-los}
\includegraphics[width=0.85\textwidth]{grafico1.png}
\fuente{Fuente: Elaboración propia}
\end{figure}

\begin{figure}[!htbp]
\centering
\caption{Grado de satisfacción de los clientes de Cafetalino}\label{fig:grado-de-satisfaccion-de-los}
\includegraphics[width=0.85\textwidth]{grafico2.png}
\fuente{Fuente: Elaboración propia}
\end{figure}

\begin{figure}[!htbp]
\centering
\caption{Ocurrencia de fallos en el servicio de Cafetalino}\label{fig:ocurrencia-de-fallos-en-el}
\includegraphics[width=0.85\textwidth]{grafico3.png}
\fuente{Fuente: Elaboración propia}
\end{figure}

\begin{figure}[!htbp]
\centering
\caption{Sentimientos del cliente de Cafetalino ante ocurrencia de fallos del servicio}\label{fig:sentimientos-del-cliente-de-cafetalino}
\includegraphics[width=0.85\textwidth]{grafico4.png}
\fuente{Fuente: Elaboración propia}
\end{figure}

\begin{figure}[!htbp]
\centering
\caption{Respuesta de Cafetalino al Cliente ante fallos del servicio}\label{fig:respuesta-de-cafetalino-al-cliente}
\includegraphics[width=0.85\textwidth]{grafico5.png}
\fuente{Fuente: Elaboración propia}
\end{figure}

\section{Diagramas de los Prototipos}\label{anexo-diagramas-de-los-prototipos}

Este anexo recoge el diagrama de arquitectura del prototipo de Visión Artificial de Bajo Costo, elaborado durante la fase de prototipado del proceso de Design Thinking y descrito en el capítulo de desarrollo conceptual. Se trata de la alternativa que resultó descartada en la selección del prototipo, por lo que su diagrama se conserva aquí como evidencia del proceso de ideación y no en el cuerpo del documento. El diagrama muestra las capas del sistema, los componentes que las integran y el flujo de información entre ellas. La arquitectura de la alternativa finalmente seleccionada, el Sistema de reabastecimiento predictivo y enrutamiento dinámico, se presenta en la Figura~\ref{fig:estructura-del-sistema-de-reabastecimiento} del apartado~\ref{arquitectura-y-pila-tecnologica}, donde se describe componente por componente.

\begin{figure}[!htbp]
\centering
\includegraphics[height=0.82\textheight,keepaspectratio]{image2.png}
\caption{Estructura de la Aplicación Visión Artificial de Bajo Costo}\label{fig:estructura-de-la-aplicacion-vision}
\fuente{Fuente: Elaboración propia}
\end{figure}

\section{Capturas de la Gestión del Proyecto en Jira}\label{anexo-capturas-de-jira}

Este anexo reúne las capturas de pantalla de la herramienta Jira con la que el equipo gestionó las historias de usuario del proyecto, descritas en el capítulo de metodología. La primera muestra la lista de tareas pendientes del proyecto con algunas de las historias de usuario consideradas; la segunda, una iteración en ejecución con sus historias clasificadas según su estado, pendiente, en ejecución, en revisión y terminadas; y la tercera, el conjunto de iteraciones planificadas para el desarrollo del sistema. En conjunto documentan el uso efectivo de la herramienta a lo largo de las cuatro iteraciones previstas.

\begin{figure}[!htbp]
\centering
\includegraphics[width=0.94\textwidth]{image4.jpeg}
\caption{Lista de tareas pendientes del proyecto (Product Backlog)}\label{fig:lista-de-tareas-pendientes-del}
\fuente{Fuente: Elaboración propia}
\end{figure}

\begin{figure}[!htbp]
\centering
\includegraphics[width=0.95\textwidth]{image5.jpeg}
\caption{Registro de tareas de una iteración en ejecución}\label{fig:registro-de-tareas-de-una}
\fuente{Fuente: Elaboración propia}
\end{figure}

\begin{figure}[!htbp]
\centering
\includegraphics[width=0.62\textwidth]{image6.jpeg}
\caption{Iteraciones del Proyecto}\label{fig:iteraciones-del-proyecto}
\fuente{Fuente: Elaboración propia}
\end{figure}


\end{document}
